%%%%%%%%%%%%%%%%%%%%%%%%%%%%%%%%%%%%
%       A free boundary problem in combustion
%                 Revised version  December  1999
%                 sent to editors December 1999
%%%%%%%%%%%%%%%%%%%%%%%%%%%%%%%%%%%%

\documentstyle[12pt,leqno]{amsart}
\vsize=27.9cm
\hsize=21cm
\setlength{\textwidth}{16cm}
\setlength{\textheight}{21.5cm}
\setlength{\oddsidemargin}{0.0cm}
\setlength{\evensidemargin}{0.0cm}
%\setlength{\topmargin}{-1.5cm}
\parskip 6pt
%%%%%%%%%%%%%%%%%%%%%%%%%%%%%%%%%
\newcommand{\ve}{\varepsilon}
\newcommand{\car}{{\cal X}}
\newcommand{\R}{{\Bbb R}}
\newcommand{\cald}{{\cal D}}
\newcommand{\bep}{{\cal B}_\ve}
\newcommand{\wep}{w^\ve}
\newcommand{\uep}{u^\ve}
\newcommand{\vep}{v^\ve}
\newcommand{\uepk}{u^{\ve_k}}
\newcommand{\vepk}{v^{\ve_k}}
\newcommand{\wepk}{w^{\ve_k}}
\newcommand{\ep}{\varepsilon}
\newcommand{\epk}{{\ve_k}}
\newcommand{\epn}{\ve_{n}}
\newcommand{\epnk}{\ve_{n_k}}
\newcommand{\uepn}{u^{\ve_{n}}}
\newcommand{\uepnk}{u^{\ve_{n_k}}}
\newcommand{\uepkn}{u^{\ve_{k_n}}}
\newcommand{\vepkn}{v^{\ve_{k_n}}}
\newcommand{\wepkn}{w^{\ve_{k_n}}}
\newcommand{\epkn}{{\ve_{k_n}}}
\newcommand{\bepk}{{\cal B}_{\epk}}
\newcommand{\uepj}{u^{\ve_j}}
\newcommand{\vepj}{v^{\ve_j}}
\newcommand{\wepj}{w^{\ve_j}}
\newcommand{\epj}{{\ve_j}}
\newcommand{\bepj}{{\cal B}_{\epj}}
\newcommand{\vepjn}{v^{\ve_{j_n}}}
\newcommand{\uepjn}{u^{\ve_{j_n}}}
\newcommand{\wepjn}{w^{\ve_{j_n}}}
\newcommand{\epjn}{{\ve_{j_n}}}
\newcommand{\vepjk}{v^{\ve_{j_k}}}
\newcommand{\uepjk}{u^{\ve_{j_k}}}
\newcommand{\wepjk}{w^{\ve_{j_k}}}
\newcommand{\epjk}{{\ve_{j_k}}}
\newcommand{\pepj}{P_{\ve_j}}
\newtheorem{defi}{Definition}[section]
\newtheorem{lema}{Lemma}[section]
\newtheorem{cor}{Corollary}[section]
\newtheorem{prop}{Proposition}[section]
\newtheorem{teo}{Theorem}[section]
%\theoremstyle{remark}
\newtheorem{rem}{Remark}[section]
%\renewcommand{\therem}{}
\newtheorem{ack}{Acknowledgements}
\renewcommand{\theack}{}
\renewcommand{\theequation}{\arabic{section}.\arabic{equation}}

\begin{document}

\title{\bf A free boundary problem in combustion theory}

\author[J. Fernandez Bonder]{Juli\'an Fern\'andez Bonder}
\author[N. Wolanski]{Noemi  Wolanski}
\thanks{Partially supported by Universidad de Buenos Aires under grant TX47,
by ANPCyT PICT No. 03-00000-00137 and CONICET PIP0660/98.
N. I. Wolanski is a member of CONICET}
\address{Departamento  de Matem\'atica, FCEyN, UBA (1428) 
Buenos Aires, Argentina.}
\email{\tt jfbonder@@dm.uba.ar \\ 
wolanski@@dm.uba.ar}

\maketitle

\begin{abstract}
In this paper we consider the following problem arising in combustion theory
$$\left\{\begin{array}{rcll}
\Delta u^{\ve} - u^{\ve}_{t} & = & v^{\ve} f_{\ve}(u^{\ve}) 
& \mbox{in } {\cal D},\\
\Delta v^{\ve} - v^{\ve}_{t} & = & v^{\ve} f_{\ve}(u^{\ve}) 
& \mbox{in } {\cal D},\\
\end{array}\right.$$
where ${\cal D}\subset \R^{N+1}$, $f_\ve(s) = \frac1{\ve^2}
f(\frac s\ve)$ with $f$ a Lipschitz continuous function with support in
$(-\infty,1]$.

Here $\vep$ is the mass fraction of some reactant, $\uep$ the rescaled
temperature of the mixture and $\ep$ is essentially the inverse of the activation energy.
This model is derived in the framework of the theory of
equidiffusional premixed flames for Lewis number 1.

We prove that, under suitable assumptions on the functions $\uep$ and $\vep$, 
we can pass to the limit ($\ve\to0$) -- the so called high activation energy
limit -- and that the limit function
$u=\lim u^{\ve}=\lim v^\ve$ is a solution of the following free boundary problem
\begin{equation}
\begin{aligned}
&\Delta u - u_t = 0\quad \mbox{in } \{ u>0 \}\\
&|\nabla u| = \sqrt{2M(x,t)}\quad \mbox{on } \partial\{ u>0 \}
\end{aligned}
\end{equation}
 in a pointwise sense at regular free boundary 
points and in a viscosity sense. Here $M(x,t)=\int_{-w_0(x,t)}^1(s+w_0(x,t))
f(s)ds$ and $-1<w_0=\lim_{\ve\to0}\frac{v^\ve-u^\ve}{\ve}$.

Since $\vep-\uep$ is a solution of the heat equation it is fully determined by
its initial-boundary datum. In particular, the free boundary condition only (but
strongly) depends
on the approximation of the initial-boundary datum.

Moreover, if ${\cal D}\cap\partial\{u>0\}$ is a Lipschitz surface, $u$ is
a classical solution to (0.1).
\end{abstract}

\begin{section}{Introduction}
\setcounter{equation}{0}
In this paper we consider the following problem arising in combustion theory

\begin{equation}
\label{1.1}
\left\{\begin{array}{rcll}
\Delta u^{\ve} - u^{\ve}_{t} & = & v^{\ve} f_{\ve}(u^{\ve}) 
& \mbox{in } {\cal D},\\
\Delta v^{\ve} - v^{\ve}_{t} & = & v^{\ve} f_{\ve}(u^{\ve}) 
& \mbox{in } {\cal D},\\
\end{array}\right.
\end{equation}
where ${\cal D}\subset\R^{N+1}$.

This model appears in combustion theory in the analysis of the propagation 
of curved flames. It is derived in the framework of the theory of 
equidiffusional premixed flames analyzed in the relevant limit of high 
activation energy for Lewis number 1. In this application, $v^{\ve}$ represents the fraction of 
some reactant (and hence it is assumed to be nonnegative), and $u^{\ve}$ its 
temperature (more precisely, $u^{\ve}=\lambda(T_f-T^{\ve})$ where $T_f$ is 
the flame temperature  and $\lambda$ is a 
normalization factor). Observe that the term $v^{\ve}f_{\ve}(u^{\ve})$ acts 
as an absorption term in the equation (\ref{1.1}). 
Since $T^{\ve} = T_f - (u^{\ve}/\lambda)$, it is in fact a reaction term for 
the temperature. In the flame model, such a term represents the effect of the 
exothermic chemical reaction and $f$ has accordingly a number of properties: 
it is a nonnegative Lipschitz continuous function which is positive in an 
interval $(-\infty, \ve)$ and vanishes otherwise (i.e., reaction occurs only 
when $T>T_f-\frac\ve\lambda$). The parameter $\ve$ is essentially the inverse of the activation 
energy of the chemical reaction. For the sake of simplicity we will
assume that $f_{\ve}(s)=\frac{1}{\ve^2}f(\frac{s}{\ve}),$ where
$f$ is a nonnegative, Lipschitz continuous function with support in 
$(-\infty,1]$.

For the derivation of the model, we cite \cite{BL}.

Here we are interested in high activation energy limits (i.e. $\ve\to 0$).
These limits, are currently the subject of active investigation, specially
 in the case
$u^{\ve}=v^{\ve}$. This is a natural assumption in the case of traveling waves.

The study of the limit as $\ep\to 0$ was proposed
in the  30's by Zeldovich and Frank-Kamenetski \cite{ZF} and has been much 
discussed in the combustion literature. 
In the case $\uep=\vep$ the reaction function  $\uep f_\ep(\uep)$  
tends to a Dirac delta, $M_0\delta(u)$ where $M_0=\int_0^1sf(s)ds$. 
In this way the reaction  
zone where $\uep f_\ep(\uep)$  acts is reduced to a surface, the flame 
front,  and a free boundary problem arises. The fact that $M_0>0$ 
ensures that a nontrivial combustion process takes place so that a 
non-empty free  boundary actually appears.


Although  the  convergence of  
the most relevant propagation modes, i.e. the  traveling  waves, was 
already discussed by Zeldovich and  Frank-Kamenetski, and an
enormous progress in this direction has been made, 
a rigorous  mathematical investigation of the 
convergence of
general solutions is still in progress. Berestycki and his
collaborators have rigorously studied the  convergence  problem for
traveling waves and, more generally in the  elliptic stationary case,
cf. \cite{BCN} and  its references. See also \cite{LW}. The study  of the 
limit in the general evolution case for the heat operator has been 
performed
in  \cite{CV} for the one  phase case (this is, with $\uep\ge0$) and in \cite{C, CLW1, CLW2} 
for the two-phase 
case, where no sign restriction on $\uep$ is made.



 In \cite{CV} the authors show that, under certain assumptions on the initial
datum and its approximations, for every sequence $\epn\to0$ there exists
a subsequence $\epnk$ and a limit  function $u=\lim u^{\ve_{n_k}}$ 
which solves
the following free boundary problem 
\begin{equation}
\label{1.1p}
\left\{\begin{aligned}
&\Delta u - u_t  =  0  \quad\ \ \ \mbox{in } {\cal D}\cap \{u>0\}\\
&|\nabla u^+|  =  \sqrt{2M_0} \quad \mbox{on } {\cal D}\cap\partial\{u>0\}
\end{aligned}\right.
\end{equation}
in a weak integral sense. Here $M_0=\int^1_0sf(s)ds$.


In  \cite{CLW1} and \cite{CLW2} the authors
show that the free boundary condition for the two phase 
case (when it is assumed that no reaction takes place if $\uep\le 0$) is 
$$|\nabla u^+|^2 - |\nabla u^-|^2 = 2M_0$$ 
and that the limit function is a solution of the
free boundary problem in a pointwise sense at
regular free boundary points when $\{ u=0 \}$ has zero ``parabolic density'' 
and in a viscosity sense in the absence of a zero phase (i.e. when 
$\{ u=0 \}^\circ\cap{\cal D}=\emptyset$)


On the other hand, in \cite{LVW} it was shown for the one phase problem
in a cylinder with Neumann boundary conditions,
that when a classical solution to the free boundary problem (\ref
{1.1p}) exists, it is the
limit of the whole family $\uep$, not just of a subsequence. Moreover, this
classical solution is the limit of $\uep$ independently of the choice of the 
approximate initial data $\uep(x,0)$. A similar result has been obtain for the
two phase case in \cite{LVW2}.

So that a natural question is if a classical solution to the free boundary
problem (\ref{1.1p}) is also the limit of $\uep$ if $(\uep,\vep)$ is the solution
to the system (\ref{1.1}) and both $\uep(x,0)$ and $\vep(x,0)$ converge 
to $u(x,0)$  but $\uep(x,0)\neq\vep(x,0)$.

Or we may ask a more elementary question: Will a sequence of uniformly
bounded solutions $(\uep,\vep)$ with $(\vep-\uep)\to0$ as $\ep\to0$ be such
that $\uep$ converges to a solution of the free boundary problem (\ref{1.1p})?
This is, will the asymptotic limit for activation energy going to infinity, in the case
in which $(\vep-\uep)\to0$ but $\uep\neq\vep$, be a solution of the same free boundary
problem as in the case in which $\uep=\vep$?

In order to understand the relation between both assertions it is important
to point out that in the case under consideration this is, when Lewis number
is 1, the function $\wep=\vep-\uep$ is a solution of the heat equation. So that
it is fully determined by its initial-boundary datum. Moreover,
the system (\ref{1.1}) may be rewritten as a single equation for $\uep$, namely
\begin{equation}
\label{reduced}
\Delta\uep-\uep_t=(\uep+\wep)f_\ep(\uep).
\end{equation}

In this paper we consider the case in which $\wep/\ep$ converges to a 
function $w_0$ (so that in particular, $\vep-\uep\to0$). In this way, at least
formally, the reaction term still converges to a delta function and a free
boundary problem appears. {\it But we prove that the free boundary condition strongly
depends on the limit function $w_0$, so that it is different for different
approximations of the initial-boundary datum of $u$}.

In fact, we prove that for every sequence $\epn\to0$ there exists a subsequence
$\epnk$ and a limit function $u=\lim \uepnk$ which is a solution of the
following free boundary problem
\begin{equation}
\label{FBP}
\left\{
\begin{aligned}
&\Delta u - u_t  =  0  \quad\ \ \ \mbox{in } {\cal D}\cap \{u>0\}\\
&|\nabla u^+|  =  \sqrt{2M(x,t)} \quad \mbox{on } {\cal D}\cap\partial\{u>0\}
\end{aligned}
\right.
\end{equation}
where $M(x,t)=\int_{-w_0(x,t)}^1\big(s+w_0(x,t)\big)f(s)ds$.

The presence of the function $w_0$ in the limit of integration gives the
necessary positive sign of the function $M(x,t)$.

In conclusion, {\it the combustion problem is very unstable in the sense that
the asymptotic limit
for activation energy going to infinity depends on order $\ep$ perturbations
of the initial-boundary data}.

In this paper we prove that the limit function $u$ is a ``viscosity''
solution to (\ref{FBP}), so that, as a consequence of our results and of
the regularity results for viscosity solutions to (\ref{FBP}) of \cite{F},
we deduce that, when the free boundary of a limit function $u$ is given by $x_1 = g(x',t)$, 
$x=(x_1,x')$ with $g$ Lipschitz continuous, $u$ is a classical solution. 



We want to stress, that because of our
assumption that $\vep-\uep\to0$ and since $\vep\ge0$, the limit function $u$ must be nonnegative, so our
result is new even in the case $u^\ve=v^\ve$.

In particular, as a consequence
of our results we see that limit functions $u$ with $u^\ve(x,0)$
 constructed as in \cite{CV}, and $v^\ve(x,0)$ small perturbations of $u^\ve
(x,0)$
 are viscosity solutions to (\ref{FBP}). In this construction, $w_0$ is any
constant such that $w_0\ge -\eta$ where $\eta>0$ is small enough.

 \bigskip

\noindent{\bf Notation, hypotheses and outline of the paper}


Throughout this paper $N$ will denote the spatial dimension and, in addition, 
the following notation will be used:

For any $x_0\in\R^N,\ t_0\in\R$ and $\tau>0$
\begin{align*}
B_{\tau}(x_0) &\equiv \{ x\in \R^N/\ |x-x_0|<\tau \},\\
B_{\tau}(x_0,t_0) &\equiv \{ (x,t)\in \R^{N+1}/\ |x-x_0|^2+|t-t_0|^2<\tau^2 \},\\
Q_{\tau}(x_0,t_0) &\equiv B_{\tau}(x_0)\times (t_0-\tau^2,t_0+\tau^2),\\
Q^-_{\tau}(x_0,t_0) &\equiv B_{\tau}(x_0)\times (t_0-\tau^2,t_0],
\end{align*}
and for any set $K\subset \R^{N+1}$

\begin{align*}
{\cal N}_{\tau}(K) &\equiv \bigcup_{(x_0,t_0)\in K} Q_{\tau}(x_0,t_0),\\
{\cal N}^-_{\tau}(K) &\equiv \bigcup_{(x_0,t_0)\in K} Q^-_{\tau}(x_0,t_0).
\end{align*}

When necessary, we will denote points in $\R^N$ by $x=(x_1,x')$, with 
$x'\in\R^{N-1}$. Also, $\langle \cdot,\cdot \rangle$ will mean the usual 
scalar product in $\R^{N}$. Given a function $v$, we will denote 
$v^+ = \max(v,0),\ v^-=\max(-v,0)$.

In addition, the symbols $\Delta$ and $\nabla$ will denote the corresponding 
operators in the space variables; the symbol $\partial_p$ will denote 
parabolic boundary.

Finally, we will say that a function $v$ is in the class $Lip_{loc}(1,1/2)$ 
in a domain ${\cal D}\subset \R^{N+1}$, if for every ${\cal D}'\subset\subset 
{\cal D}$, there exists a constant $L=L({\cal D}')$ such that
$$|v(x,t) - v(y,s)| \le L (|x-y| + |t-s|^{1/2})$$
for every $(x,t),(y,s)\in {\cal D}'$. If the constant $L$ does not depend on 
the set ${\cal D}'$, we will say that $v\in Lip(1,1/2)$ in ${\cal D}$.

\bigskip

For the existence of a limit function for a subsequence
$\uepnk$ we only need the weaker condition
that
for every compact $K\subset {\cal N}_{\tau}^-(K)\subset{\cal D}$, 
\begin{equation}
\label{1.2}
\| v^{\ve} - u^{\ve} \|_{L^{\infty}({\cal N}_{\tau}^-(K))} = \mbox{O}(\ve).
\end{equation}
Then, we have (see \cite{LSU})
\begin{equation}
\label{1.3}
\| v^{\ve} - u^{\ve} \|_{C^{2,1}(K)} = \mbox{O}(\ve).
\end{equation}

Under this assumption, we are able to apply the results of \cite{C} and get 
the uniform Lipschitz estimates needed to pass to the limit in (\ref{1.1}). 
 This is done in \S 2 where
we also prove some technical lemmas that are 
used throughout the paper.

\medskip

In \S 3 and \S 4 we assume that $u^{\ve}\to 0$ in $\{ u=0 \}$ fast enough. 
This is an essential condition that was already considered in \cite{CV}.
This assumption is a natural one in applications, roughly speaking it means 
that the mixture temperature reaches the flame temperature only if some 
combustion is taking place.

 We also assume that
there exists $\lim_{\ve\to0}(v^\ve- u^\ve)/\ve=:w_0$ and, as a consequence
of the hypothesis that $\uep\to0$ in $\{u=0\}$ fast enough, we
show that necessarily $w_0>-1$ in $\{u\equiv0\}^\circ$.
So that, in \S 3 and \S 4 we assume that for
every $K\subset \cal N_\tau^-(K)\subset
\cald$ compact
\begin{equation}
\label{ochica}
\frac{ v^{\ve} - u^{\ve}}{\ve}\to w_0 \qquad\mbox{uniformly in }N_\tau^-(K).
\end{equation}
Thus,
\begin{equation}
\label{ochica2}
\big\| \frac{v^{\ve} - u^{\ve}}{\ve}-w_0\big \|_{C^{2,1}(K)} \to 0.
\end{equation}
And, for the sake of simplicity, we  assume that $w_0>-1$ in  $\cald$.

\medskip

In \S 3, we show that the limit function $u$ is a solution to the free
boundary problem (\ref{FBP}) 
in a pointwise sense. 

\medskip

Finally, in \S 4 we prove
that the limit function $u$ is in fact a viscosity solution of the free
boundary problem 
(\ref{FBP}) under a nondegeneracy assumption on the limit function $u$. We also
prove some results that give the necessary nondegeneracy of $u$.

Our presentation is of a local nature, so that our hypotheses are stated in terms of the
solution $(\uep,\vep)$.  As can be seen in the example
treated in Corollary \ref{cor4}  it is possible to deduce our hypotheses on $(\uep,\vep)$ from 
conditions on its initial-boundary datum.

\end{section}

\bigskip
%%%%%%%%%%%%%%%%%%%%%%%%%%%%%%%%%%%%%%%%%%%%%%%%%%%%%%%%%%%%
%%%%			SECTION 2			%%%%
%%%%%%%%%%%%%%%%%%%%%%%%%%%%%%%%%%%%%%%%%%%%%%%%%%%%%%%%%%%%


\begin{section}{Uniform estimates}
\setcounter{equation}{0}

In this section we consider a family $u^{\ve}, v^{\ve}$ of solutions to
(\ref{1.1}) in a domain $\cal D$ which are uniformly bounded in $L^{\infty}$ 
norm in $\cal D$. We show that the functions $u^{\ve}, v^{\ve}$ are locally 
uniformly bounded in the seminorm $Lip(1,1/2)$.
Then, we get further local uniform estimates and pass to the limit as 
$\ve\to 0$. We also show that the limit function $u$ is a solution to the 
free boundary problem (\ref{FBP}) in a very weak sense. Finally, we prove an approximation 
lemma that will be used throughout the rest of the paper and some lemmas
concerning particular limit functions. Also, we prove a proposition that
justifies the hypothesis we make in the following sections.

For convenience, let us define the following function
\begin{equation}
\label{w}
w^{\ve}(x,t) = v^{\ve}(x,t) - u^{\ve}(x,t),
\end{equation}
then, $w^\ve$ is a caloric function with $\| w^\ve\|_{C^{2,1}(K)}=\mbox{O}(\ve)$ for 
every compact set $K\subset {\cal D}$.

We begin with a proposition (which is a consequence of \cite{C}) that gives 
us the uniform control on the gradients of solutions of (\ref{1.1}).

\begin{prop}
\label{regx}
Let $(u^{\ve},v^{\ve})$ be solutions of $(\ref{1.1})$ such that 
$\| u^{\ve} \|_{\infty}\leq {\cal A}$, $v^{\ve}\ge0$ and verify $(\ref{1.2})$. 
Let $K\subset {\cal D}$ compact and $\tau>0$ such that 
${\cal N}_{\tau}^{-}(K)\subset {\cal D}$. Then, there exists 
$L=L(\tau,{\cal A})$ such that
$$|\nabla u^{\ve}(x,t)|\leq L,\quad |\nabla v^{\ve}(x,t)|\leq L.$$
\end{prop}

\begin{pf}
Let us start by making the following observation
$$u^{\ve}=v^{\ve}-w^{\ve}\ge -w^{\ve}\ge -C\ve.$$

Then, let $z^{\ve}=\frac{1}{C+1}(u^{\ve}+C\ve)$ and we define, for 
$(x_0,t_0)\in K$
$$z^{\ve}_{\tau}(x,t)=\frac{1}{\tau}z^{\ve}(x_0+\tau x,t_0+\tau^2 t).$$
In $B_1(0)\times [-1,0]$, $z^{\ve}_{\tau}$ verifies  (with $B\ge \|f\|_\infty$)
\begin{align*}
0\le\Delta z^{\ve}_{\tau}-\frac{\partial z^{\ve}_{\tau}}{\partial t}&\le
\frac{\tau}{C+1}(C\ve+|u^{\ve}|)\frac{1}{\ve^2}f(\frac{u^{\ve}}{\ve})\\
&\le B{\tau}\frac{1}{\ve}\car_{[-C\ve,\ve]}(u^{\ve})=
\frac{B}{\ve/\tau}\car_{[0,\ve/\tau]}(z^{\ve}_{\tau}).
\end{align*}
On the other hand
$$|z^{\ve}_{\tau}(x,t)|\le\frac{ |u^{\ve}(x,t)|+C}{\tau(1+C)}\le
\frac{1}{\tau}\frac{{\cal A}+C}{1+C}.$$
Therefore, by Theorem 2 of \cite{C} it follows that
$$|\nabla z^{\ve}_{\tau}(x,t)|\le \bar L=\bar L(\tau,{\cal A})\qquad 
\mbox{in } B_{1/2}(0)\times(-1/2,0].$$
In particular,
\begin{align*}
|\nabla u^{\ve}(x_0,t_0)| &=(C+1)|\nabla z^{\ve}(x_0,t_0)|=
(C+1)|\nabla z^{\ve}_{\tau}(0,0)|\le (C+1)\bar L,\\
|\nabla v^{\ve}(x_0,t_0)| &\le|\nabla u^{\ve}(x_0,t_0)|+
|\nabla w^{\ve}(x_0,t_0)|\le (C+1)\bar L+C.
\end{align*}
The proof is finished\end{pf}



As is usual in parabolic theory, Lipschitz regularity in space, gives 
H\"older 1/2 regularity in time. We follow here ideas in \cite{CLW1} and 
\cite{CV}.

\begin{prop}
\label{regt}
Let $(u^{\ve}, v^{\ve})$ be solutions of $(\ref{1.1})$ such that 
$\| u^{\ve} \|_{\infty}\leq {\cal A}$, $v^{\ve}\ge0$, and verify $(\ref{1.2})$. 
Let $K\subset {\cal D}$ compact and $\tau>0$ such that 
${\cal N}_{\tau}(K)\subset {\cal D}$. Then there exists $C=C(\tau,{\cal A})$ 
such that
$$| u^{\ve}(x,t+\Delta t) - u^{\ve}(x,t)|\leq C |\Delta t|^{1/2},\quad
| v^{\ve}(x,t+\Delta t) - v^{\ve}(x,t)|\leq C |\Delta t|^{1/2}, $$
for every $(x,t), (x,t+\Delta t)\in K$.
\end{prop}

\begin{pf}
As in Proposition \ref{regx} we define $z^{\ve}=\frac{1}{C+1}(u^{\ve}+C\ve)$ and
$$z^{\ve}_{\lambda}(x,t)=\frac{1}{\lambda}z^{\ve}(x_0+\lambda x,
t_0+\lambda^2 t),$$
for $0<\lambda<\tau$ and $(x_0,t_0)\in K$.

By a simple computation we get, as in Prop. 2.1
$$0\le \Delta z^{\ve}_{\lambda} - \frac{\partial z^{\ve}_{\lambda}}{\partial t}
\le\frac{B}{\ve/\lambda}\car_{[0,\ve/\lambda]}(z^{\ve}_{\lambda}).$$
Now, $z^{\ve}_{\lambda}\ge0$, and in $\{ z^{\ve}_{\lambda}>1 \}$ we have
$$|\Delta z^{\ve}_{\lambda} - \frac{\partial z^{\ve}_{\lambda}}{\partial t}|
\left\{\begin{array}{cc}
		\le B & \mbox{ if } \ve/\lambda\ge1\\
		= 0   & \mbox{ if } \ve/\lambda<1.\\
	\end{array}\right.$$
Moreover, we have that
$$|\nabla z_{\lambda}^{\ve}(x,t)|=
\frac1{C+1}|\nabla u^{\ve}(x_0+\lambda x,t_0+\lambda^2 t)|\le \bar L$$
in $B_{\tau/\lambda}(0)\times[0,\tau^2/\lambda^2]$.
Then, by Proposition 2.2 of \cite{CLW1}, we have
$$|z^{\ve}_{\lambda}(0,t)-z^{\ve}_{\lambda}(0,0)|\le C(\bar L)\qquad
\forall\ 0\le t\le\frac{1}{4N+B}$$
which, in terms of $u^{\ve}$, is
$$|u^{\ve}(x_0,t_0+\lambda^2 t)-u^{\ve}(x_0,t_0)|\le C(\bar L)\lambda.$$
In particular
$$|u^{\ve}(x_0,t_0+\frac{\lambda^2}{4N+B})-u^{\ve}(x_0,t_0)|\le C(\bar L)\lambda.$$

Let $(x_0,t_0+\Delta t)\in K$. If $0<\Delta t<\frac{\tau^2}{4N+B}$, we take
$\lambda=\Delta t^{1/2}\sqrt{4N+B}<\tau$ to get
$$|u^{\ve}(x_0,t_0+\Delta t)-u^{\ve}(x_0,t_0)|\le C(\bar L)\sqrt{4N+B}
\Delta t^{1/2}.$$
If $\Delta t\ge\frac{\tau^2}{4N+B}$, we have
$$|u^{\ve}(x_0,t_0+\Delta t)-u^{\ve}(x_0,t_0)|\le 2{\cal A}\le
\frac{2{\cal A}}{\tau}\sqrt{4N+B}\Delta t^{1/2}.$$

The analogous inequality for $v^{\ve}$ is an immediate consequence of 
(\ref{1.3}).\end{pf}

\medskip

\begin{rem} Under the hypothesis of the previous propositions, 
we have that 
$$u^{\ve}\in Lip_{\mbox{loc}}(1,1/2).$$
\end{rem}

\medskip

\begin{prop}
\label{lim}
Let $(u^{\ve}, v^{\ve})$ be solutions of $(\ref{1.1})$ such that 
$\| u^{\ve} \|_{\infty}\leq {\cal A}$, $ v^{\ve}\ge0$ and verify
$(\ref{1.2})$. Then, for every sequence
$\ve_{n}\to 0$, there exists $\ve_{n'}\to 0$ a subsequence and 
$u\in Lip_{\mbox{loc}}(1,1/2)$ such that
\begin{enumerate}
\item $u^{\ve_{n'}}\to u$ uniformly on compacts subsets of ${\cal D}$.
\item $\nabla u^{\ve_{n'}}\to \nabla u$ in $L^{2}_{\mbox{loc}}$.
\item $\frac{\partial u^{\ve_{n'}}}{\partial t}\to 
\frac{\partial u}{\partial t}$ weakly in $L^{2}_{\mbox{loc}}$.
\item $\Delta u - \frac{\partial u}{\partial t} = 0$ in 
$\{ u>0 \}$
\item For every compact $K\subset {\cal D}$, exists $C_{K}>0$ such that
$$\left\| \frac{\partial u^{\ve}}{\partial t} \right\|_{L^{2}(K)}\leq C_{K}$$
for every $\ve>0$.
\end{enumerate}
\end{prop}

\begin{pf}
The proof is similar to lemma 3.1 of \cite{CLW1}.

Let $K\subset{\cal D}$ be a compact set, and $\tau>0$ such that
${\cal N}_{3\tau}(K)\subset{\cal D}$.
Let $L=L(K)$ such that
$$| u^{\ve}(x,t)-u^{\ve}(y,s) |\le L\left( |x-y|+|t-s|^{1/2}\right),$$
where $(x,t), (y,s)\in {\cal N}_{\tau}(K)$.

Then, by Arzela-Ascoli's theorem, there exists $\ve_{n^{\prime}}\to 0$ and
$u\in Lip(1,1/2)$ in ${\cal N}_{\tau}(K)$ such that
$u^{\ve_{n^{\prime}}}\to u$ uniformly in ${\cal N}_{\tau}(K)$.
By a standard diagonal argument, (1) follows.

Let us now find uniform bounds for $\frac{\partial u^{\ve}}{\partial t}$ in
$L^2_{\mbox{loc}}({\cal D})$. In fact, $u^{\ve}$ verifies
$$\Delta u^{\ve} - \frac{\partial u^{\ve}}{\partial t} = 
v^{\ve}f_{\ve}(u^{\ve}).$$

Now, let $(x_0,t_0)\in K$ and let us multiply the equation by 
$u^{\ve}_{t}\psi^2$ where $\psi\ge0$, 
$\psi=\psi(x)\in C^{\infty}_{c}(B_{\tau}(x_0))$, $\psi\equiv 1$ in 
$B_{\tau/2}(x_0)$. Then, integrating by parts, we get

\begin{align*}
\iint_{Q_{\tau}(x_0,t_0)}(u^{\ve}_{t})^2\psi^2 dxdt &+
\frac12 \iint_{Q_{\tau}(x_0,t_0)} (|\nabla u^{\ve}|^2)_t \psi^2 dxdt +
2 \iint_{Q_{\tau}(x_0,t_0)} \nabla u^{\ve} u^{\ve}_t \psi \nabla \psi dxdt\\
& = -\iint_{Q_{\tau}(x_0,t_0)} v^{\ve}f_{\ve}(u^{\ve})u^{\ve}_t \psi^2 dxdt.
\end{align*}
Now we use Young's inequality to obtain

\begin{align*}
&\frac12 \iint_{Q_{\tau}(x_0,t_0)}(u^{\ve}_{t})^2\psi^2 dxdt +
\frac12 \int_{B_{\tau}(x_0)} |\nabla u^{\ve}(x_0,t_0+\tau^2)|^2 \psi^2 dx \le\\
&\frac12 \int_{B_{\tau}(x_0)} |\nabla u^{\ve}(x_0,t_0-\tau^2)|^2 \psi^2 dx 
- \iint_{Q_{\tau}(x_0,t_0)} v^{\ve}f_{\ve}(u^{\ve})u^{\ve}_t \psi^2 dxdt\\ 
&+C\iint_{Q_{\tau}(x_0,t_0)}|\nabla u^{\ve}|^2|\nabla \psi|^2 dxdt.
\end{align*}
Then, by Proposition \ref{regx}
\begin{align*}
&\int_{B_{\tau/2}(x_0)}\int_{t_0-\tau^2}^{t_0+\tau^2}(u^{\ve}_{t})^2 dxdt\le
\int_{B_{\tau}(x_0)} |\nabla u^{\ve}(x_0,t_0-\tau^2)|^2 \psi^2 dx \\
&+2\left|\iint_{Q_{\tau}(x_0,t_0)} v^{\ve}f_{\ve}(u^{\ve})u^{\ve}_t \psi^2 
dxdt\right|
+ C\iint_{Q_{\tau}(x_0,t_0)}|\nabla u^{\ve}|^2|\nabla \psi|^2 dxdt\\
&\le C(\tau) +2 \left|\iint_{Q_{\tau}(x_0,t_0)} v^{\ve}f_{\ve}(u^{\ve})
u^{\ve}_t \psi^2 dxdt\right|.
\end{align*}
Hence, it only remains to get bounds on
$$\iint_{Q_{\tau}}\psi^{2} u^{\ve}_{t} v^{\ve}f_{\ve}(u^{\ve}) dx dt = I.$$

Let 
$${\cal G}_{\ve}(u,x,t) = \int_{0}^{u}(w^{\ve}(x,t)+s)f_{\ve}(s) ds,$$
then
$$\frac{\partial}{\partial t}\left({\cal G}_{\ve}(u^{\ve},x,t)\right)=
\frac{\partial u^{\ve}}{\partial t} v^{\ve}f_{\ve}(u^{\ve}) +
\frac{\partial {\cal G}_{\ve}}{\partial t}(u^{\ve},x,t),$$
so that we get
$$I=\iint_{Q_{\tau}}\psi^{2}\frac{\partial}{\partial t}\left(
{\cal G}_{\ve}(u^{\ve},x,t)\right) dx dt
- \iint_{Q_{\tau}}\psi^{2}\frac{\partial {\cal G}_{\ve}}
{\partial t}(u^{\ve},x,t) dx dt = A - B.$$
Let us first get bounds on $A$:
\begin{align*}
A&=\int_{t_{0}-\tau^{2}}^{t_{0}+\tau^{2}}\int_{B_{\tau}(x_{0})}
\psi^{2}\frac{\partial}{\partial t}\left({\cal G}_{\ve}(u^{\ve},x,t)\right) dx dt =
\int_{B_{\tau}(x_{0})}\psi^{2}\left[\int_{t_{0}-\tau^{2}}^{t_{0}+\tau^{2}}
\frac{\partial}{\partial t}\left({\cal G}_{\ve}(u^{\ve},x,t)\right) dt \right] dx\\
&=\int_{B_{\tau}(x_{0})}\psi^{2}\left[ 
{\cal G}_{\ve}(u^{\ve}(x,t_0+\tau^2),x,t_{0}+\tau^{2})-
{\cal G}_{\ve}(u^{\ve}(x,t_0-\tau^2),x,t_{0}-\tau^{2})\right] dx.
\end{align*}
Since $u^{\ve}\ge -C\ve$, $f_{\ve}(s)=0$ if $s\ge\ve$ and 
$|w^{\ve}|=\mbox{O}(\ve)$, we have 
$$|{\cal G}_{\ve}(u^{\ve},x,t)|\leq C\ve\int_{-C\ve}^{\ve}
f_{\ve}(s) ds+\int_{-C\ve}^{\ve}sf_{\ve}(s)ds\le C,$$ 
so that
$$|A|\leq C(\tau).$$

It only remains to get bounds on $B$. For that purpose, let us first make the 
following observation:
$$\left|\frac{\partial {\cal G}_{\ve}}{\partial t}(u^{\ve},x,t)\right|=
\left|\frac{\partial w^{\ve}}{\partial t}(x,t)
\int_{0}^{u^{\ve}}f_{\ve}(s) ds\right|
\leq \frac{C}{\ve}\left|\frac{\partial w^{\ve}}{\partial t}(x,t)\right|.$$
By (\ref{1.3}), 
$$\left| \frac{\partial w^{\ve}}{\partial t}\right|\le C\ve 
\quad  \mbox{ for } (x,t)\in{\cal N}_{\tau}(K).$$
Therefore, using the fact that $0\leq \psi \leq 1$, we get
$$B\leq \frac{C}{\ve}\iint_{Q_{\tau}}\left|\frac{\partial w^{\ve}}
{\partial t}(x,t)\right| dx dt  \le \frac C\ve|Q_\tau|
\left| \frac{\partial w^{\ve}}{\partial t}\right|\leq  C(K,\tau).$$
Thus, 
$$\int_{B_{\tau/2}(x_0)}\int_{t_0-\tau^2}^{t_0+\tau^2} (u^{\ve}_t)^2 dxdt
\le C,$$ 
with $C$ independent of $\ve$ and $(x_0,t_0)\in K$. Now, as $K$ is compact, 
$$\iint_{K}(u^{\ve}_t)^2 dxdt\le C,$$
so that, for a subsequence, $\frac{\partial}{\partial t}u^{\ve_{n^\prime}}
\to \frac{\partial}{\partial t}u$ weakly in $L^2(K)$ and by a standard 
diagonal argument, (3) follows.

Let us see that $u$ is a solution of the heat equation in $\{ u>0\}$ .
In fact, from the fact that $u^{\ve}\to u$ uniformly on compact subsets of 
${\cal D}$, we deduce that every point $(x_0,t_0)\in\{ u>0\}$ has a 
neighborhood $V$ such that $u^{\ve}(x,t)\ge\lambda>0$ for some $\lambda>0$. 
Therefore, for $\ve<\lambda$, $f_{\ve}(u^{\ve}(x,t))=0$ in $V$. Thus 
$u^{\ve}$ is caloric in $V$ for every $\ve<\lambda$, and then, the same fact 
holds for $u$.

Let us finally analyze the convergence of the gradients. We already know that 
\linebreak $\|\nabla u^{\ve}\|_{L^{\infty}({\cal N}_{\tau}(K))}\le L$. So we can assume 
that $\nabla u^{\ve}\to \nabla u$ weakly in $L^{2}({\cal N}_{\tau}(K))$. 
In particular
$$\iint_{{\cal N}_{\tau}(K)}\phi |\nabla u|^2\le 
\liminf_{\ve\to0}\iint_{{\cal N}_{\tau}(K)}\phi |\nabla u^{\ve}|^2,$$
for every nonnegative $\phi\in L^{\infty}({\cal D})$.

We follow here ideas from \cite{BCN} and \cite{CV} in order to prove that we 
have strong convergence.

Since $\Delta u - u_t = 0$ in $\{ u>0 \}$, if we take $\delta>0$ and multiply 
this equation by $(u-\delta)^+ \psi(x)$ with $\psi\in L^{\infty}({\cal D})$ 
and nonnegative, we get after integration by parts in $Q_{\tau}(x_0,t_0)$, 
\begin{align*}
&\iint_{\{ u>\delta\}}|\nabla u|^2 \psi =
-\iint_{\{ u>\delta\}}u\nabla u\nabla \psi +
\delta\iint_{\{ u>\delta\}}\nabla u\nabla\psi \\
&- \frac12 \int_{\{ u>\delta\}} (u-\delta)^2(x,t_0+\tau^2)\psi(x)
+ \frac12 \int_{\{ u>\delta\}} (u-\delta)^2(x,t_0-\tau^2)\psi(x).
\end{align*}
Now, letting $\delta\to 0$, we get
\begin{align*}
\iint_{\{ u>0\}}|\nabla u|^2\psi=&
-\iint_{\{ u>0\}} u\nabla u\nabla\psi -\\
&-\frac12\int_{\{ u>0\}} u^2(x,t_0+\tau^2)\psi(x)
+ \frac12\int_{\{ u>0\}} u^2(x,t_0-\tau^2)\psi(x).
\end{align*}

On the other hand, since $\psi\ge0$, $f_{\ve}\ge0$ and $u^{\ve}\ge-C\ve$, 
multiplying (\ref{1.1}) by $(u^{\ve}+C\ve)\psi$ and integrating by parts 
we get
\begin{align*}
&\iint_{Q_{\tau}(x_0,t_0)}|\nabla u^{\ve}|^2\psi\le
-\iint_{Q_{\tau}(x_0,t_0)} u^{\ve}\nabla u^{\ve}\nabla\psi -
C\ve\iint_{Q_{\tau}(x_0,t_0)} \nabla u^{\ve}\nabla\psi \\
&- \frac12\int_{B_{\tau}(x_0)} (u^{\ve}+C\ve)^2(x,t_0+\tau^2)\psi(x) +
\frac12\int_{B_{\tau}(x_0)} (u^{\ve}+C\ve)^2(x,t_0-\tau^2)\psi(x).
\end{align*}
Thus, 
$$\limsup_{\ve\to0}\iint_{Q_{\tau}(x_0,t_0)}|\nabla u^{\ve}|^2\psi\le
\iint_{Q_{\tau}(x_0,t_0)}|\nabla u|^2\psi,$$
so that
$$\|\psi^{1/2}\nabla u^{\ve}\|_{L^2(Q_{\tau}(x_0,t_0))}\to
\|\psi^{1/2}\nabla u\|_{L^2(Q_{\tau}(x_0,t_0))}.$$
Since, in addition, 
$$\psi^{1/2}\nabla u^{\ve}\to\psi^{1/2}\nabla u
\ \mbox{ weakly in } L^2(Q_{\tau}(x_0,t_0)),$$
it follows that
$$\psi^{1/2}\nabla u^{\ve}\to\psi^{1/2}\nabla u
\ \mbox{ in } L^2(Q_{\tau}(x_0,t_0)).$$
Therefore, as $\psi\equiv 1$ in $B_{\tau/2}(x_0)$, 
$$\nabla u^{\ve}\to\nabla u\ \mbox{ in } L^2(Q_{\tau/2}(x_0,t_0))$$
and since $K$ is compact, this implies that
$$\nabla u^{\ve}\to\nabla u\ \mbox{ in } L^2(K).$$
By the same standard diagonal argument used before, the assertion of the 
theorem follows. 
\end{pf}



Next we show that the limit function $u$ is a solution of the free
boundary problem in a very 
weak sense.


\begin{prop}
Let $(u^{\ve_{j}}, v^{\ve_{j}})$ be a family of solutions of $(\ref{1.1})$ in a 
domain ${\cal D}\subseteq \R^{N+1}$ such that $u^{\ve_{j}}\to u$ 
uniformly on compact subsets of ${\cal D}$, $v^{\ve_{j}}\ge0$ and verify 
$(\ref{1.2})$. Then, there exists a locally finite measure $\mu$ supported on 
the free boundary ${\cal D}\cap \partial \{ u>0 \}$ such that
$v^{\ve_j}f_{\ve_{j}}(u^{\ve_{j}})\to \mu$ weakly in ${\cal D}$ and 
therefore
$$\Delta u - \frac{\partial u}{\partial t} = \mu\ \mbox{ in } {\cal D},$$
that is $\forall \ \phi \in C^{\infty}_{c}({\cal D})$
$$\iint_{\cal D} (u\phi_{t} - \nabla u \nabla \phi) dx dt = 
\iint_{\cal D} \phi\ d\mu.$$
\end{prop}

\begin{pf}
The proof follows as that of proposition 3.1 of \cite{CLW1}
\end{pf}


Now we state an approximation lemma that will be used throughout the rest of 
the paper.

\begin{lema} 
\label{lim2}
Let $(u^{\ve_{j}}, v^{\ve_{j}})$ be a family of solutions of $(\ref{1.1})$ in a 
domain ${\cal D}\subseteq \R^{N+1}$ such that $u^{\ve_{j}}\to u$ 
uniformly on compact subsets of ${\cal D}$, $v^{\ve_{j}}\ge0$ and verify 
$(\ref{1.2})$. Let $(x_0,t_0)\in{\cal D}\cap\partial\{ u>0 \}$ and let 
$(x_n,t_n)\in{\cal D}\cap\partial\{ u>0 \}$ be such that 
$(x_n,t_n)\to (x_0,t_0)$ as $n\to \infty$. 
Let $\lambda_n\to 0,\ u_{\lambda_n}(x,t)=\frac{1}{\lambda_n}
u(x_n+\lambda_n x, t_n+\lambda_n^2 t)$ and $(u^{\ve_j})_{\lambda_n}(x,t)=
\frac{1}{\lambda_n} u^{\ve_j}(x_n+\lambda_n x, t_n+\lambda_n^2 t)$.
Assume that $u_{\lambda_n}\to U$ as $n\to\infty$ uniformly on compact sets 
of $R^{N+1}$. Then, there exists $j(n)\to\infty$ such that for every 
$j_n\ge j(n)$ there holds that $\frac{\ve_{j_n}}{\lambda_n}\to 0$ and
\begin{enumerate}
\item $(u^{\ve_{j_n}})_{\lambda_n}\to U$ uniformly on compact sets of 
$R^{N+1}$,
\item $\nabla(u^{\ve_{j_n}})_{\lambda_n}\to \nabla U$ in $L^2_{loc}(R^{N+1})$, 
\item $\frac{\partial}{\partial t}(u^{\ve_{j_n}})_{\lambda_n}\to
\frac{\partial}{\partial t}U$ weakly in $L^2(R^{N+1})$.

Also, we deduce that
\item $\nabla u_{\lambda_n}\to \nabla U$ in $L^2(R^{N+1})$,
\item $\frac{\partial}{\partial t}u_{\lambda_n}\to 
\frac{\partial}{\partial t}U$ weakly in $L^2(R^{N+1})$.
\end{enumerate}
\end{lema}

\begin{pf}
The proof is a straightforward adaptation of lemma 3.2 of \cite{CLW1}
\end{pf}

Now we state some lemmas on special limits of solutions to $P_{\ve}$ that
will be used throughout the paper.

\begin{lema}
\label{global1}
Let $(\uepj,\vepj)$ be a solution to $(1.1)$ in a domain $\cald\subset\R^{N+1}$
such that $\vepj\ge0$,  and verify $(\ref{ochica})$ in $\cald$ with
$w_0=\mbox{constant}$. 
Let $(x_0,t_0)\in\cald$ and assume that $\uepj$ converges
to $u=\alpha(x-x_0)_1^+$ uniformly on compact subsets of $\cald$, with $\alpha\in\R$
and $\epj\to0$. Then,
\begin{equation}
0\le\alpha\le\sqrt{2M}.
\end{equation}
where $M=\int_{-w_0}^1 (s+w_0)f(s)\,ds$.
\end{lema}

\begin{pf}
The proof is an adaptation of Proposition 5.2 of \cite{CLW1}.

Without loss of generality we may assume that $(x_0,t_0)=(0,0)$.

First we see that necessarily $\alpha\ge0$ since $u$ is subcaloric in $\cald$
and $u(0,0)=0$. If  $\alpha=0$ there is nothing to prove. So let us assume
that $\alpha>0$. 

Let $\psi\in C_c^\infty(\cald)$. Multiplying $(1.1)$ by $u^\ve_{x_1}\psi$
and integrating by parts we get
\begin{equation}
\begin{aligned}
\label{equepx1}
\iint_{\cald}\uepj_t\uepj_{x_1}\psi=&\frac12\iint_{\cald}|\nabla\uepj|^2\psi_{x_1}
-\iint_{\cald}\uepj_{x_1}\nabla\uepj\nabla\psi\\
&\iint_{\cald}\bepj(\uepj,x,t)\psi_{x_1}
+\iint_{\cald}\wepj_{x_1}\big(\int_{-w_0}^{\uepj}f_{\epj}(s)ds\big)\psi,
\end{aligned}
\end{equation}
where $\bep(u,x,t)=\int_{-w_0\ve}^u(s+\wep)f_\ve(s)ds$.

In order to pass to the limit in $(\ref{equepx1})$ we observe that, by Proposition
$\ref{lim}$
$$
(\uepj)_t\to 0\qquad\mbox{weakly in }L^2_{loc}(\cald),
$$
$$
\nabla\uepj\to\alpha\car_{\{x_1>0\}} e_1\qquad\mbox{in }L^2_{loc}(\cald).
$$
On the other hand,
$$
\frac{\nabla\wepj}{\epj} \to 0\qquad\mbox{uniformly on compact subsets of }\cald.
$$
Therefore, in order to pass to the limit in $(\ref{equepx1})$ we only need
to analyze the limit of $\bepj(\uepj,x,t)$. On one hand, it is easy to see
that
\begin{equation}
\label{limbepj}
\bepj(\uepj(x,t),x,t)\to M
\end{equation}
for every $(x,t)$ such that $x_1>0$. In fact,
$$
\bepj(\uepj,x,t)=\int_{-w_0}^{\frac{\uepj}{\epj}}\big(s+\frac{\wepj}{\epj}\big)
f(s)ds=
\int_{-w_0}^1\big(s+\frac{\wepj}{\epj}\big)f(s)ds
$$
if $j$ is large enough. Since $|\bepj(\uepj,x,t)|\le C$ there holds that
(\ref{limbepj}) holds in $L^1_{loc}(\{x_1\ge0\}$.

On the other hand, there exists $\bar M(x,t)\in L^\infty(\cald)$ such that
$\bepj(\uepj,x,t)\to \bar M(x,t)$ weakly in $L^2_{loc}(\cald)$. Clearly,
$\bar M(x,t)=M$ in $\{x_1>0\}$. Let us see that
$\bar M(x,t)=\bar M(t)$ in $\{x_1<0\}$. In fact,
\begin{align*}
\nabla(\bepj(\uepj(x,t),x,t))&=\frac{\partial\bepj}{\partial u}(\uepj,x,t)
\nabla\uepj+\nabla\bepj(\uepj,x,t)\\
&=(\uepj+\wepj)f_{\epj}(\uepj)\nabla\uepj+\nabla\wepj\int_{-w_0\epj}^{\uepj}
f_{\epj}(s)ds\\
&= \vepj f_{\epj}(\uepj)\nabla\uepj+ \frac{\nabla\wepj}{\epj}\int_{-w_0}
^{\frac{\uepj}{\epj}}f(s)ds.
\end{align*}
Since $\vepj f_{\epj}(\uepj)\to 0$ in $L^1_{loc}(\{x_1<0\})$, $\nabla\uepj$
is uniformly bounded in $L^\infty(\cald')$ if $\cald'\subset\subset\cald$ and
$\frac{\nabla\wepj}{\epj}\to 0$ uniformly on compact subsets of $\cald$, there
holds that
$$
\nabla(\bepj(\uepj(x,t),x,t))\to 0\qquad\mbox{in }L^1_{loc}(\{x_1<0\}).
$$
So that, passing to the limit in (\ref{equepx1}) we get
$$
\frac{\alpha^2}2\iint_{\{x_1>0\}}\psi_{x_1}=M\iint_{\{x_1>0\}}\psi_{x_1}+
\int_{\{x_1<0\}}\bar M(t)\psi_{x_1}.
$$
Thus, integrating in the variable $x_1$ we get
$$
\int_{\{x_1=0\}}\big(\frac{\alpha^2}2-M+\bar M(t)\big)\psi=0.
$$
Since $\psi$ is arbitrary, we conclude that
$$
\frac{\alpha^2}2-M+\bar M(t)=0.
$$
Finally, we notice that $\bar M(t)\ge0$. In fact,
\begin{align*}
\bepj(\uepj,x,t)=\int_{-\frac{\wepj}{\epj}}^{\frac{\uepj}{\epj}}&\big(s+\frac{\wepj}
{\epj}\big)f(s)ds +\int_{-w_0}^{-\frac{\wepj}{\epj}}\big(s+\frac{\wepj}{\epj}\big)
f(s)ds\\
&\ge\int_{-w_0}^{-\frac{\wepj}{\epj}}\big(s+\frac{\wepj}{\epj}\big) f(s)ds\to0
\end{align*}
since  $ \frac{\wepj}{\epj}\to w_0$ uniformly on compact subsets of $\cald$.

            Thus,
            $$
            \alpha=\sqrt{2(M-\bar M(t))}\le\sqrt{2M}
            $$
            and the proof is complete.\end{pf}


\begin{lema}
\label{global2}
Let $(\uepj,\vepj)$ be a solution to $(1.1)$ in a domain $\cald\subset\R^{N+1}$
such that $\vepj\ge0$ and verify $(\ref{ochica})$ with $w_0=\mbox{constant}$
in $\cald$.
Let $(x_0,t_0)\in\cald$ and assume that $\uepj$ converges
to $u=\alpha(x-x_0)_1^++\bar\alpha(x-x_0)_1^-$ uniformly on compact subsets of 
$\cald$, with $\alpha,\bar\alpha>0$ and $\epj\to0$. Then,
\begin{equation}
\bar\alpha=\alpha\le\sqrt{2M}
\end{equation}
where $M=\int_{-w_0}^1 (s+w_0)f(s)\,ds$.
\end{lema}


\begin{pf}
We argue in a similar way as in Proposition 5.3 of \cite{CLW1}.

We will denote $Q_r = Q_r(0,0)$. Without loss of generality we will assume 
that $(x_0,t_0)=(0,0)$ and that $Q_2\subset\subset \cal D$.

As before, $u^{\ve}$ satisfies 
\begin{align*}
\iint_{\cal D} u^{\ve}_t u^{\ve}_{x_1}\psi =&
\frac12\iint_{\cal D}|\nabla u^{\ve}|^2 \psi_{x_1} - 
\iint_{\cal D} u^{\ve}_{x_1}\nabla u^{\ve}\nabla\psi\\
& + \iint_{\cal D}{{\cal B}}_{\ve}(u^{\ve},x,t)\psi_{x_1} + 
\iint_{\cal D}w^{\ve}_{x_1}\left(\int_0^{u^{\ve}} f_{\ve}(s) ds\right)\psi.
\end{align*}

We want to pass to the limit. By Proposition \ref{lim} and the fact that 
$u^{\ve_j}$ converge to $\alpha x_1^+ + \bar{\alpha} x_1^-$ we have 
that
$$u^{\ve_j}_t\to 0  \qquad \mbox{weakly in } L^2_{loc}(\cal D),$$
$$\nabla u^{\ve_j}\to \alpha \car_{\{x_1>0\}}e_1 - 
\bar{\alpha} \car_{\{x_1<0\}} e_1
\quad \mbox{in } L^2_{loc}(\cal D).$$

Clearly, as $\alpha,\bar{\alpha}>0$, ${\cal B}(u^{\ve_j},x,t)\to M$ in 
$L^1_{loc}(\cal D)$.

So, passing to the limit in the latter equation for the subsequence $\ve_j$, 
we get 
$$ - \frac{\alpha^2}{2}\iint_{\{x_1>0\}}\psi_{x_1} - 
\frac{\overline{\alpha}^2}{2}\iint_{\{x_1<0\}}\psi_{x_1}+
M\iint \psi_{x_1} = 0.$$
Integrating in the $x_1$ variable, we conclude that 
$$\alpha = \bar{\alpha}.$$

Next, we will assume that $\alpha>\sqrt{2M}$ and arrive at a contradiction.

First, let us consider  $z^{\ve_j}$, defined in 
$Q_2$, the solution to 

\begin{equation}
\label{eqz}
\Delta z^{\ve_j} - z^{\ve_j}_t = \left(\beta_{\ve_j}(z^{\ve_j}) +
W_{\ve_j} f_{\ve_j}(z^{\ve_j})\right)\rho_{\ve_j}(z^{\ve_j}/\epj)\quad\mbox{ in }
Q_2
\end{equation}
with  boundary conditions
$$
z^{\ve_j} = u - b^{\ve_j} \quad\mbox{ on } \partial_p Q_2
$$
where $\beta_{\ve}(s) = sf_{\ve}(s)$, $W_{\ve} = \sup_{Q_2} w^{\ve}$, 
$b^{\ve_j} = \sup_{Q_2}|u^{\ve_j} - u|$ and $\rho_{\ve_j}$ is a smooth cutoff 
function with support in $[-(w_0+2C_{\ve_j}),2]$ and $\rho_{\ve_j}\equiv1$ in 
$[-(w_0+C_{\ve_j}),1]$ (Here $C_{\epj}\to0^+$ is such that $\big|\wepj/\epj
-w_0\big|\le C_{\epj}$ in $Q_2$
so that $\uepj/\epj\ge-(w_0+C_{\epj})$ in $Q_2$).


Observe that $z^{\ve_j}(x_1,x',t) = z^{\ve_j}(-x_1,x',t)$ in $Q_2$.

It is easy to see that the proofs of Propositions \ref{regx} and \ref{regt}
can be adapted to $z^{\ve_j}$
so that, for a subsequence, that we still call $\epj$, there holds that
$z^{\ve_j}\to z$ uniformly on compact sets of $Q_2$.
We will show that $z = u$.

First,
\begin{align*}
\Delta\uepj-\uepj_t &=\vepj f_{\epj}(\uepj)=(\uepj+\wepj) f_{\epj}(\uepj) \le 
\beta_{\epj}(\uepj)+W_{\epj} f_{\epj}(\uepj)  \\
&=\big(\beta_{\epj}(\uepj)+W_{\epj} f_{\epj}(\uepj)\big) \rho_{\epj}(\uepj/\epj)
 \quad\qquad\mbox{ in } Q_2.
\end{align*}


From the fact that $z^{\ve_j}\le u^{\ve_j}$ on $\partial_p Q_2$,
we deduce that $z^{\ve_j}\le u^{\ve_j}$ in $Q_2$ and therefore $z\le u$.

In order to see that $u\le z$, we consider $a^{\ve_j}\in C^2(\R)$ such that 
$$\begin{array}{l}
a^{\ve_j}_{ss} = \left(\beta(a^{\ve_j}) + \frac{W_{\ve_j}}{\ve_j} f(a^{\ve_j})
\right)\rho_{\ve_j}(a^{\ve_j}),\ s\in\R\\
a^{\ve_j}(0) = 1,\ a^{\ve_j}_s(0) = \alpha
\end{array}$$

Integrating the equation we get, for every $s\in\R$, that
$$0<\gamma - \kappa_{\ve_j}\le a^{\ve_j}_s(s)\le \alpha$$
where $\frac12 \gamma^2 \equiv \frac12 \alpha^2 - M >0$ and 
$\kappa_{\ve_j}\to 0$ when $j\to\infty$.

It follows that there exists $\overline{s}_{\ve_j}<0$ such that
$$a^{\ve_j}(s) = \left\{\begin{array}{ll}
1+\alpha s & s\ge0\\
(\gamma - \kappa_{\ve_j})(s-\overline{s}_{\ve_j}) & s\le\overline{s}_{\ve_j}
\end{array}\right.$$
and it is easy to see that $\overline{s}_{\ve_j}$ are uniformly bounded by 
below and moreover, there exists $\overline{s}<0$ such that
$\overline{s}_{\ve_j}\to\overline{s}$.

Now let
$$\tilde{a}^{\ve_j}(x) = \ve_j a^{\ve_j}\big(\frac{x_1}{\ve_j} - 
\frac{b^{\ve_j}}{(\gamma - \kappa_{\ve_j})\ve_j} + \overline{s}_{\ve_j}\big).$$

Using that $\tilde{a}^{\ve_j}(0,x',t) = -b^{\ve_j}$ and the bounds on 
$a^{\ve_j}_s$, we deduce that
$$\tilde{a}^{\ve_j}\le u - b^{\ve_j} \mbox{ in } Q_2.$$

Now, since $\tilde{a}^{\ve_j}\le z^{\ve_j}$ on $\partial_p Q_2$, and
$\tilde a^{\epj}$ is a one dimensional stationary solution to (\ref{eqz}), we have that
$\tilde{a}^{\ve_j}\le z^{\ve_j}$ in $Q_2$. Since $\tilde{a}^{\ve_j}\to u$ 
uniformly on compact subsets of $\{ x_1>0 \}$, we deduce that 
$u\le z$ in $Q_2\cap \{ x_1>0 \}$.

Finally, we notice that $z^{\ve_j}(x_1,x',t) = z^{\ve_j}(-x_1,x',t)$, so we 
conclude that $u\le z$ in $Q_2$.

\medskip

Now, the proof of the Lemma follows as in \cite{CLW1} where it is shown, for 
the case in which $w^{\ve}\equiv 0$, that if $z^{\epj}\to \alpha x_1^+
+\alpha x_1^-$  uniformly on compact subsets of $Q_2$, where $z^{\epj}$ are
solutions to  (\ref{eqz}) which are symmetric in the $x_1$ variable, 
and $\alpha>0$, there holds that
$\alpha\le \sqrt{2M}$. So that we arrive at a contradiction since we have
assumed that $\alpha>\sqrt{2M}$. Here we use again that $\frac{\wepj}{\epj}
\to w_0$ and  $\mbox{liminf}\,\bepj(\uepj,x,t)\ge0$. The proof is finished.
\end{pf}

\begin{lema}
\label{global3}
Let $(\uepj,\vepj)$ be a solution to $(1.1)$ in a domain ${\cald}_j$ such that
$\vepj\ge0$ and
satisfies $(\ref{ochica})$ in ${\cald}_j$ with $w_0=\mbox{constant}$.
Here ${\cald}_j$ is such
that ${\cald}_j\subset\cald_{j+1}$ and $\cup_j{\cald}_j=\R^{N+1}$. Let us assume
that $\uepj\to U$ uniformly on compact subsets of ${\R}^{N+1}$ as $j\to\infty$
and $\epj\to0$, with $U\ge0$,  $U\in Lip(1,1/2)$  and $\partial\{U>0\}\neq
\emptyset$.  Then,
\begin{equation}
|\nabla U|\le \sqrt{2M}\qquad\mbox{in }\R^{N+1}
\end{equation}
with $M=\int_{-w_0}^1(s+w_0)f(s)ds$.
\end{lema}

\begin{pf} The proof is similar to that of Theorem 6.2 in \cite{CLW1}. Here
we use Lemmas \ref{global1} and \ref{global2} instead of Propositions 5.2
and 5.3 in \cite{CLW1}.
\end{pf}



\end{section}

\bigskip
%%%%%%%%%%%%%%%%%%%%%%%%%%%%%%%%%%%%%%%%%%%%%%%%%%%%%%%%%%%%%%%%%%%%
%%%%			SECTION 3                               %%%%
%%%%%%%%%%%%%%%%%%%%%%%%%%%%%%%%%%%%%%%%%%%%%%%%%%%%%%%%%%%%%%%%%%%%

\begin{section}{The free boundary condition}
\setcounter{equation}{0}

In this section, we find the free boundary condition for the limit problem 
and we show that the limit function $u$ is a solution to the free boundary
problem (\ref{FBP}) 
in a pointwise sense, under the assumption that the free boundary admits an 
inward spatial normal in a parabolic measure theoretic sense (Definition 
\ref{def1}). 

Throughout this section we will assume that
$(\ref{ochica})$ holds
and that
for every
$K\subset\{u\equiv0\}^\circ$
compact there exists $0<\eta<1$ and $\ve_0>0$ such that, for $\ve<\ve_0$
\begin{equation}
\label{umenordelta}
\frac{u^\ve}\ve\le \eta\qquad\mbox{in }K.
\end{equation}

This assumption is a natural one in applications, roughly speaking it means
that the mixture temperature reaches the flame temperature only if some
combustion is taking place.

As a consequence there holds that
\begin{equation*}
w_0=\lim_{\ep\to0}\frac{\vep-\uep}\ep\ge-\limsup_{\ep\to0}\frac{\uep}\ep\ge-\eta
>-1\quad\mbox{in }K.
\end{equation*}
So that, for the sake of simplicity we will assume from now on that $w_0>-1$
in $\cald$.

We start this section with a lemma that is the essential ingredient in the 
subsequent proofs.

\begin{lema} 
\label{Bepsilon}
Let $(u^{\ve_k},v^{\ve_k})$ be a solution to $(1.1)$ in a domain $\cald\subset
\R^{N+1}$ such that $v^{\ve_k}\ge0$ and $(\ref{ochica})$ and
$(\ref{umenordelta})$ are satisfied with $w_0>-1$. Let
$u=\lim u^{\ve_k}$, with $\ve_k\to0$, and $\bepk(u,x,t)=\int_{-w_0\epk}^{u}
\big(s+{\wepk})f_{\epk}(s)ds$. Then,
$$
{{\cal B}}_{\ve_k}(u^{\ve_k},x,t)\to M(x,t)\car_{ \{ u>0 \} },\quad
\mbox{in } L^1_{loc}(\cald).
$$
where $M(x,t)=\int_{-w_0(x,t)}^1(s+w_0(x,t))f(s)ds$.
\end{lema}

\begin{pf}
First, let us observe that 
\begin{align*}
\int_{-w_0\ve}^{\ve}(w^{\ve}+s)f_{\ve}(s)ds&=\int_{-w_0\ve}^{\ve}
(w^{\ve}+s)\frac{1}{\ve^2}f\left(\frac{s}{\ve}\right)ds\\
&=\int_{-w_0}^1\big(\frac{w^\ve}{\ve}+s\big)f(s)ds.
\end{align*}
Therefore, 
$$\lim_{\epk\to0}\int_{-w_0\epk}^{\epk}(\wepk+s)f_{\epk}(s)ds=M(x,t).$$

\medskip

Let us now see that ${\cal B}_{\ve_k}(u^{\ve_k},x,t)\to M(x,t)$
uniformly on compact subsets of $\{ u>0 \}$.

Let $K\subset\subset \{ u>0 \}$, then there exists $\lambda>0$ and $\ve_0$
such that $u^{\ve_k}(x,t)>\lambda\ \ \forall \ve_k<\ve_0\,,\,(x,t)\in K$.
Thus, we have

\begin{align*}
\lim_{k\to\infty}{\cal B}_{\ve_k}(u^{\ve_k}(x,t),x,t)&=
\lim_{k\to\infty}\int_{-w_0\epk}^{u^{\ve_k}(x,t)}(w^{\ve}+s)f_{\ve_k}(s)ds\\
&=\lim_{k\to\infty}\int_{-w_0\epk}^{\ve_k}(w^{\ve}+s)f_{\ve_k}(s)ds=M(x,t).
\end{align*}

 
Since $|\bepk(\uepk,x,t)|\le C$ on every compact subset of $\cald$,
there holds, for a subsequence that we still call $\epk$ that
$$
\bepk(\uepk,x,t)\to \bar M(x,t)\qquad\mbox{weakly in }L^2_{loc}(\cald).
$$
Clearly, $\bar M(x,t)=M(x,t)$ in $\{u>0\}$. Let us see that $\bar M(x,t)=0$
in $\{u\equiv0\}^\circ$. In fact, let $K$ be a compact subset of $\{u\equiv0\}
^\circ$.
For every $\ve_1,\ve_2>0$ there holds that, for $k$ large enough
\begin{align*}
&|\{(x,t)\in K\,/\,\ve_1<\bepk(\uepk,x,t)<M(x,t)-\ve_2\}|\\
&\quad\le |\{(x,t)\in K\,/\,\frac{\uepk}{\epk}(x,t)>-w_0(x,t)\,,\,\frac
{\ve_1}2<\int_{-w_0}^{\frac{\uepk}{\epk}}(s+w_0)f(s)ds<M-\frac{\ve_2}2\}|\\
&\quad\le |\{(x,t)\in K\,/\, -w_0(x,t)+\mu\le\frac{\uepk}{\epk}\le 1-\mu\}|\\
&\quad\le |\{(x,t)\in K\,/\,\frac{\vepk}{\epk}\ge\frac\mu2\,,\,\frac{\uepk}
{\epk}\le 1-\mu\}|\\
&\quad\le|\{(x,t)\in K\,/\, \vepk f_{\epk}(\uepk)\ge \frac{C_\mu}{2\epk}\}|.
\end{align*}
Since $\vepk f_{\epk}(\uepk)\to 0$ as measures in $K$ and $\vepk\ge0$,
$f_{\epk}\ge0$ there holds that
$$
\vepk f_{\epk}(\uepk)\to 0 \qquad\mbox{in }L^1(K).
$$
Therefore,
$$
|\{ (x,t)\in K\,/\,\ve_1<\bepk(\uepk,x,t)<M(x,t)-\ve_2\}|\to 0.
$$
On the other hand,
let $1>\eta>\sup_K(-w_0)$ be the constant in (\ref{umenordelta})  in $K$, there holds that
\begin{align*}
&\bepk(\uepk,x,t)= \int_{-\frac{\wepk}{\epk}}^{\frac{\uepk}{\epk}}\big(
s+\frac{\wepk}{\epk}\big)f(s)ds+
\int_{-w_0}^{-\frac{\wepk}{\epk}}\big(s+\frac{\wepk}{\epk}\big)f(s)ds\\
&\qquad\le  \int_{-\frac{\wepk}{\epk}}^{\eta}\big(s+\frac{\wepk}{\epk}\big)
f(s)ds+
\int_{-w_0}^{-\frac{\wepk}{\epk}}\big(s+\frac{\wepk}{\epk}\big)f(s)ds\\
&\qquad=\int_{-w_0}^\eta\big(s+\frac{\wepk}{\epk}\big)f(s)ds\to
\int_{-w_0}^\eta(s+w_0)f(s)ds<M(x,t)
\end{align*}
since $-\frac{\wepk}{\epk}\le \frac{\uepk}{\epk}\le \eta$ in $K$.
Therefore,
$$
\limsup\bepk(\uepk,x,t)\le\int_{-w_0}^{\eta}(s+w_0)f(s)ds<M(x,t).
$$
 So that, for $\ve_2>0$ small we get
$$
|\{(x,t)\in K\,/\,\ve_1<\bepk(\uepk,x,t)\}|=|\{(x,t)\in K\,/\,
\ve_1<\bepk(\uepk,x,t)<M-\ve_2\}|\to 0.
$$

Let us now see that $\bar M(x,t)=0$ in $K$. As in Lemma \ref{global1} we see that
$\bar M(x,t)\ge0$. Now assume that for some $\ve_1>0$ we have $|\{\bar M(x,t)>\ve_1\}|>0$.
Then, there exists $m$ such that $|\{\bar M(x,t)>\ve_1+\frac1m\}|:=|A_m|>0$.

Now,
$$
\int_{A_m}\bepk(\uepk,x,t)\to\int_{A_m}\bar M(x,t)>\big(\ve_1+\frac1m\big)|A_m|
$$
but,
\begin{align*}
\int_{A_m}\bepk(\uepk,x,t)&
=\int_{A_m\cap\{\bepk(\uepk,x,t)>\ve_1\}}\bepk(\uepk,x,t)\\
&+\int_{A_m\cap\{\bepk(\uepk,x,t)\le \ve_1\}}\bepk(\uepk,x,t).
\end{align*}
Since the first term in the right hand side goes to zero and the second is
bounded by $\ve_1|A_m|$, we get a contradiction.

The proof is finished.
\end{pf}

Let us give the definition of regular point.

\begin{defi}
\label{def1}
We say that $\nu$ is the interior unit spatial normal to the free boundary 
$\partial\{ u>0 \}$ at a point $(x_0,t_0)\in\partial\{ u>0 \}$ in the 
parabolic measure theoretic sense, if $\nu\in \R^N,\ |\nu|=1$ and
$$\lim_{r\to 0}\frac{1}{r^{N+2}}\iint_{Q_r(x_0,t_0)}
|\car_{\{ u>0 \}}-\car_{\{ (x,t)/\ \langle x-x_0,\nu \rangle>0\}}| dxdt=0.$$
\end{defi}

\begin{defi}
\label{def2}
We say that $(x_0,t_0)$ is a regular point of $\partial\{ u>0 \}$ if there 
exists an interior unit spatial normal to $\partial\{ u>0 \}$ at $(x_0,t_0)$ in 
the parabolic measure theoretic sense.
\end{defi}

We can  now prove the main result of this section.



\begin{teo}
\label{puntual}
Let $(u^{\ve_j},v^{\ve_j})$ be a family of uniformly bounded solutions of $(\ref{1.1})$ in a domain
${\cal D}\subset \R^{N+1}$ such that $u^{\ve_j}\to u$ uniformly on 
compact subset of ${\cal D}$, $v^{\ve_j}\ge0$ and verify $(\ref{ochica})$
and $(\ref{umenordelta})$, with $w_0>-1$.  If $(x_0,t_0)$ is 
a regular point of $\cald\cap\partial\{ u>0 \}$, then  
$u$ has the asymptotic development
$$u(x,t) = \alpha\langle x-x_0,\nu \rangle^+ + 
\mbox{o}(|x-x_0|+|t-t_0|^{1/2})$$
with $\alpha = \sqrt{2M(x_0,t_0)}$, where $M=\int_{-w_0}^1 (s+w_0)f(s)\,ds$.
Here $\nu$ is the interior unit 
spatial normal
to the free boundary at $(x_0,t_0)$ in the parabolic measure theoretic sense.
\end{teo}

\begin{pf}
We assume, without loss of generality, that $(x_0,t_0)=(0,0)$ and $\nu=e_1=(1,0,...,0)$.

Let $\psi\in C^{\infty}_c({\cal D})$. We proceed as in Lemma \ref{global1}.
Let us multiply the equation for
$u^\ve$ by $u_{x_1}^{\ve}\psi$ and integrate by parts. We have
\begin{align*}
\iint_{\cal D}u^{\ve}_t u^{\ve}_{x_1}\psi = &\frac12 \iint_{\cal D} 
|\nabla u^{\ve}|^2 \psi_{x_1} - 
\iint_{\cal D}u^{\ve}_{x_1}\nabla u^{\ve} \nabla \psi \\
&+\iint_{\cal D} {\cal B}_{\ve}(u^{\ve},x,t)\psi_{x_1} + 
\iint_{\cald}\left(\frac{w^\ve}\ve-w_0\right)f(-w_0)(w_0)_{x_1}\psi\\
&+\iint_{\cal D} \frac{w^{\ve}_{x_1}}\ve \left(\int_{-w_0}^{\frac{u^{\ve}}\ve}
f(s) ds\right)\psi
\end{align*}

Since
$$
\bepj(\uepj,x,t)=\int_{-w_0}^{\frac{\uepj}{\epj}}(s+w_0)f(s)ds +\int_{-w_0}^
{\frac{\uepj}{\epj}}\big(\frac{\wepj}{\epj}-w_0\big)f(s)ds,
$$
and $\bepj(\uepj,x,t)\to 0$ weakly in $L^1(K)$ for every $K\subset\{u\equiv0\}^\circ$
compact, there holds that
$$
F_{\epj}(x,t):=\int_{-w_0}^{\frac{\uepj}{\epj}}(s+w_0)f(s)ds\to0\qquad\mbox{weakly in }L^1(K).
$$
Since $F_{\epj}$ is nonnegative there holds that
$$
F_{\epj}\to0\qquad\mbox{in }L^1(K).
$$
So that, for a subsequence that we still call $\epj$ there holds that
$$
F_{\epj}\to0\qquad a.e.\ K.
$$
Thus,  
$$
\frac{\uepj}{\epj}\to -w_0\qquad a. e.\ K.
$$
Therefore,
\begin{equation}
\label{integralf}
\int_{-w_0}^{\frac{\uepj}{\epj}}f(s)ds\to\left(\int_{-w_0}^1f(s)ds\right)\car_
{\{u>0\}}\qquad a.e.\ \cald.
\end{equation}



By using Proposition \ref{lim}, Lemma \ref{Bepsilon} and (\ref{integralf})
we can pass to
the limit (for the sequence $\ve_j\to 0$) in the latter equation and get
\begin{equation}
\begin{aligned}
\label{u1}
&\iint_{\cal D} u_t u_{x_1}\psi = \frac12 \iint_{\cal D} 
|\nabla u|^2 \psi_{x_1} - \iint_{\cal D} u_{x_1}\nabla u \nabla \psi + 
\iint_{\{ u>0 \}}M(x,t) \psi_{x_1}\\
&\qquad+\iint_{\{u>0\}} {(w_0)}_{x_1} \left(\int_{-w_0}^1 f(s) ds\right)\psi
\end{aligned}
\end{equation}
for every $\psi\in C^{\infty}_c({\cal D})$.

Now, let $\psi^{\lambda}(x,t)=\lambda\psi(\frac{x-x_0}{\lambda},
\frac{t-t_0}{\lambda^2})$. Replacing $\psi$ by $\psi^\lambda$ in (\ref{u1}) 
and changing variables, we get for $u_{\lambda}(x,t)=
\frac{1}{\lambda}u(x_0+\lambda x,t_0+\lambda^2 t)$,
\begin{equation}
\begin{aligned}
\label{ulambda}
&\iint (u_{\lambda})_t (u_{\lambda})_{x_1}\psi = \frac12 
\iint |\nabla u_{\lambda}|^2 \psi_{x_1} - 
\iint (u_{\lambda})_{x_1}\nabla u_{\lambda} \nabla \psi\\
&\qquad +
\iint_{\{ u_{\lambda}>0 \}}M(\lambda x,\lambda^2 t) \psi_{x_1}
+\iint_{\{u>0\}} {(w_0)}_{x_1} \left(\int_{-w_0}^1 f(s) ds\right)\psi^\lambda.
\end{aligned}
\end{equation}

Let $r>0$ be such that $Q_r(x_0,t_0)\subset\subset{\cal D}$. We have that 
$u_{\lambda}\in Lip(1,1/2)$ in $Q_{r/\lambda}(0,0)$ uniformly in $\lambda$, 
and $u_{\lambda}(0,0)=0$. Therefore, for every $\lambda_n\to 0$, there 
exists a subsequence $\lambda_{n'}\to 0$ and a function $U\in Lip(1,1/2)$ in 
$\R^{N+1}$ such that $u_{\lambda_{n'}}\to U$ uniformly on compact sets of 
$\R^{N+1}$.

By our assumption on $(x_0,t_0)$, we can easily see that for every $k>0$
\begin{equation}
\label{ulposit}
|\{ u_{\lambda}>0 \}\cap \{ x_1<0 \}\cap Q_k(0,0)|\to 0
\mbox{ as } \lambda\to 0,
\end{equation}
and
\begin{equation}
\label{ulcero}
|\{ u_{\lambda}=0 \}\cap \{ x_1>0 \}\cap Q_k(0,0)|\to 0
\mbox{ as } \lambda\to 0.
\end{equation}

Now, using lemma \ref{lim2} and the fact that $\psi^\lambda
\to0$ uniformly in $\cald$ and $\mbox{supp}\,\psi^\lambda\subset\mbox{supp}\,
\psi$, we can pass to the limit in (\ref{ulambda}) and
get
\begin{equation}
\begin{aligned}
\iint_{\{ x_1>0 \}} U_t U_{x_1}\psi =&\\
 \frac12
\iint_{\{ x_1>0 \}} |\nabla U|^2& \psi_{x_1} - 
\iint_{\{ x_1>0 \}} U_{x_1}\nabla U \nabla \psi + 
M(0,0)\iint_{\{ x_1>0 \}} \psi_{x_1}.
\end{aligned}
\label{equgrande}
\end{equation}


Our aim is to prove that $U=\alpha x_1^+$. First, by (\ref{ulposit}) and
(\ref{ulcero}), we deduce that $U=0$ in $\{x_1<0\}$. On the other hand,
$U$ is a solution to the heat equation in $\{U>0\}\subset\{x_1>0\}$. By
Corollary A.1 in \cite{CLW1}, for every $\bar x'\in \R^{N-1}$, $\bar t\in \R$ there
exists $\alpha\ge0$ such that
$$
U(x,t)=\alpha x_1^++o(|(x_1,x')-(0,\bar x')|+|t-\bar t|^{1/2})\ \ \mbox{in }
\{x_1>0\}\cap\{t<\bar t\}.
$$
Replacing  the test function $\psi$ by $\Phi^\lambda(x,t)= \lambda \Phi
(\frac{x_1}\lambda,\frac{x'-\bar x'}\lambda,\frac{t-\bar t}{\lambda^2})$
with $\Phi\in C^\infty_c(\{t<0\})$ and proceeding as above we get
\begin{equation}
\label{eqalpha}
-\frac{\alpha^2}2\iint_{\{x_1>0\}}\Phi_{x_1}+M(0,0)\iint_{\{x_1>0\}} \Phi_{x_1}=0.
\end{equation}
In order to pass to the limit for a sequence $\lambda_n\to0$ we have used
Lemma \ref{lim2}. (See \cite{CLW2}, Theorem 3.1 for the details).

Thus, $\alpha=\sqrt{2M(0,0)}$.

In order to see that $U=\alpha x_1^+$ we use Lemma \ref{global3}. In fact,
by  Lemma \ref{lim2} there exists a sequence $j_n\to\infty$ such that
$$
u^{\delta_n}:=\frac 1{\lambda_n}\uepjn(\lambda_n x,\lambda_n^2 t)\to U(x,t)
$$
uniformly on compact subsets of $\R^{N+1}$. We recall that $(u^{\delta_n},v^
{\delta_n})$ is a solution to (1.1) with $\ve$ replaced by $\delta_n$.
Moreover,
$$
\frac{w^{\delta_n}}{\delta_n}=\frac{\wepjn(\lambda_n x,\lambda_n^2 t)}{\epjn}
\to w_0(0,0)
$$
uniformly on compact sets of $\R^{N+1}$.

Moreover, $U\ge0$ and $\partial\{U>0\}\neq\emptyset$. By Lemma \ref{global3}
we have that $|\nabla U|\le \alpha(=\sqrt{2M(0,0)}$. Since $U\equiv0$
in $\{x_1=0\}$ we deduce that
$$
U\le\alpha x_1\qquad\mbox{in }\{x_1>0\}.
$$
By Hopf's Principle, we deduce that
$$
U=\alpha x_1 \qquad\mbox{in }\{x_1>0\}.
$$

The theorem is proved.
\end{pf}

\medskip

\begin{rem}
It is clear from the proof that the result is still true if we replace
condition $(\ref{umenordelta})$ by the following property: $\frac{\uepj}{\epj}
\to -w_0\ \ a.e.\ \ \{u\equiv0\}^\circ$.
\end{rem}

\end{section}

\bigskip

\begin{section}{Viscosity solutions}
\setcounter{equation}{0}

In this section we prove that, under suitable assumptions, the limit function $u$ is a viscosity solution 
of the free boundary problem (\ref{FBP}).

For the sake of completeness, we state here the definition of viscosity 
solution that was introduced in \cite{CLW2} for the two phase case of this 
problem when $w_0=0$.


\begin{defi}
Let $Q$ be a cylinder in $\R^N\times (0,T)$ and let $v\in C(\overline{Q})$. 
Then $v$ is called a classical subsolution (supersolution) of $(\ref{FBP})$
in $Q$ if $v\ge0$ 
and
\begin{enumerate}
\item $\Delta v - v_t \ge 0\ (\le 0)\qquad \mbox{in } \Omega^+ \equiv 
Q\cap\{ v>0 \}$.
\item $v\in C^1(\overline{\Omega^+})$.
\item For any $(x,t)\in \partial\Omega^+\cap Q$, $\nabla v(x,t)\ne 0$, and
$$|\nabla v(x,t)|\ge \sqrt{2M}\quad (\le \sqrt{2M}).$$
\end{enumerate}

We say that $v$ is a classical solution in $Q$ if it is both a classical 
subsolution and a classical supersolution.
\end{defi}


\begin{defi}
Let $u$ be a continuous nonnegative function in $Q$; $u$ is called a 
viscosity subsolution (supersolution) of $(\ref{FBP})$ in $Q$ if, for every subcylinder 
$Q'\subset\subset Q$ and for every classical supersolution (subsolution) 
$v$ in $Q'$, 
$$\begin{array}{ll}
u\le v \ \mbox{on }\partial_p Q'  & (u\ge v\quad \mbox{on } \partial_p Q')\ \ \mbox{ and }\\
v>0 \mbox{ on } \overline{\{ u>0 \}}\cap \partial_p Q' &
(u>0 \mbox{ on } \overline{\{ v>0 \}}\cap \partial_p Q')\\
\end{array}$$
implies that $u\le v\ (u\ge v)$ in $Q'$.

The function $u$ is called a viscosity solution if it is both a viscosity 
subsolution and a viscosity supersolution.
\end{defi}

\medskip

\begin{defi}
Let $u$ be a continuous nonnegative function in $\cal D$ and let 
$(x_0,t_0)\in\partial\{ u>0 \}\cap \cal D$. We say that $(x_0,t_0)$ is a regular 
point from the nonpositive side, if there exists a regular nonnegative
function $v$ in
$\cal D$ such that $v>u$ in $\{u>0\}$ for $t<t_0$ and $v(x_0,t_0)=u(x_0,t_0)$.
\end{defi}

Finally we need the following definition on nondegeneracy.

\begin{defi}
Let $u$ be a continuous nonnegative function in $\cald$. Let $(x_0,t_0)\in\cald$
be such that $u(x_0,t_0)=0$. We say that $u$ does not degenerate at $(x_0,t_0)$
if there exist $r_0>0$ and $C>0$ such that
$$
\sup_{\partial_p Q_r^-(x_0,t_0)}u\ge C\,r\qquad\mbox{for }0<r\le r_0.
$$
\end{defi}

We now prove that, under suitable assumptions on the limit function $u$, 
there holds that $u$ is a viscosity solution to the free boundary problem.

\begin{teo} 
\label{viscosa}
Let $u=\lim u^{\ve_k}$, where $(u^{\ve_k},v^{\ve_k})$ are uniformly
bounded solutions to $(1.1)$ with   $v^{\ve_k}\ge0$, satisfying
$(\ref{ochica})$ 
in $\cal D$, with $w_0>-1$,  and such that $u^{\ve_k}$ either satisfies $(\ref{umenordelta})$
of $\uepk_t\le 0$ in $\cald$.

 If $u^+$ does not degenerate at every point of the free 
boundary which is regular from the nonpositive side, then $u$ is a viscosity 
solution of $(\ref{FBP})$.
\end{teo}

\begin{pf}

By Proposition \ref{lim} and Lemma \ref{lim2}, Theorem 4.1 of \cite{CLW2}
can be stated for our system, thus $u$ is a viscosity supersolution.

In order to see that it is a viscosity  subsolution, let $v$ be a classical
supersolution such that
$$u\le v \mbox{ in } \partial_p Q \mbox{ and } v>0 \mbox{ in }
\overline{\{u>0\}}\cap\partial_pQ$$
we want to see that $u\le v$ in $Q$.

\medskip

If not, we define
$$t_0 = \sup\{ 0<s<T:\ v>0 \mbox{ in } \overline{\{ u>0 \}}\cap Q\cap
\{ 0\le t<s \}\}$$

From the definition of $t_0$, it follows that $t_0>0$ and, from our
hypotheses we deduce that $v\ge u$ in
$Q\cap\{ 0\le t<t_0 \}$. In addition, there exists a sequence $(x(s),t(s))\to
(x_0,t_0)\in \overline{Q}$ such that $v(x(s),t(s)) = 0$, $(x(s),t(s))\in
\overline{\{ u>0 \}}\cap Q$. Clearly, $u(x_0,t_0) = v(x_0,t_0) = 0$ and 
$(x_0,t_0)\in\partial\{ u>0 \}\cap Q$. If $(x_0,t_0)\in \{ v = 0 \}^\circ$ 
then, for $\tau$ small we have $u\le v = 0$ in $B_{\tau}(x_0,t_0)\cap
\{ t<t_0 \}$ and therefore, $u\equiv 0$ there, which contradicts our 
hypothesis. Thus
$$v\ge u \mbox{ in } Q\cap \{ 0\le t\le t_0 \},$$
$$(x_0,t_0)\in \partial\{ u>0 \}\cap \partial\{ v>0 \}\cap Q.$$ 

We may assume, without loss of generality, that $(x_0,t_0)=(0,0)$ and
$Q_1(0,0)=Q_1\subset Q$ (consider instead of $u$ the function
$\frac{1}{\lambda_0}u(x_0+\lambda_0 x,t_0+\lambda_0^2 t)$ for certain $\lambda_0>0$
small, and analogously with $v$). Let us take
$$v_{\lambda}(x,t)=\frac{1}{\lambda}v(\lambda x,\lambda^2 t),\quad 
u_{\lambda}(x,t)=\frac{1}{\lambda}u(\lambda x,\lambda^2 t).$$

It is easy to see that there exists a sequence $\lambda_n\to 0$ and functions 
$u_0,\ v_0$ such that $v_{\lambda}\to v_0,\ u_{\lambda}\to u_0$.

Since $v$ is regular, we have that $v_0(x,t)=\beta x_1^+$ with 
$0\le\beta\le\sqrt{2M(0,0)}$ (for some system of coordinates).

Let us see that also $u_0(x,t)=\alpha x_1^+$ for some $\alpha\ge0$, 

We may think that in $Q_1$,  $\partial\{v>0\}$ is the graph of some function 
$\psi(x',t)=x_1,\ x=(x_1,x')$ with $\psi\in Lip(1,1/2)$, where $\psi(0,0)=0$ 
and $\{ v>0 \} = \{ x_1>\psi(x',t) \}$.

Hence, we have that
$$|\psi(x',t)|\le C\left(|x'| + |t|^{1/2}\right).$$

Let  ${\cal R}=\left\{(x,t)\in Q_1:\ x_1<-C\left(|x'| + 
|t|^{1/2}\right)\right\}$. Then ${\cal R}\cap\{ v>0 \}=\emptyset$ and let 
$w$ be the caloric function in ${\cal O}=Q_1^-\setminus {\cal R}$ with $w=0$ 
in $\partial_p{\cal R}$ and $w=L\ge\|u\|_{\infty}$ in the rest of
$\partial_p{\cal O}$.

Since $u$ is globally subcaloric and $u\le w$ on $\partial_p{\cal O}$, 
then $u < w$ in ${\cal O}$.

Now, $w-u$ is supercaloric in ${\cal O}$, $w-u>0$ in the interior and 
$w-u=0$ at $(0,0)$, then, by lemma A.1 of \cite{CLW1}, we have that 
$w-u=\delta x_1^+ + o(|x|+|t|^{1/2})$ and, since by the same lemma, $w$ has 
an asymptotic development at $(0,0)$, 
$$u(x,t)=\alpha x_1^+ + o(|x|+|t|^{1/2}),\quad \mbox{with } \alpha\ge0.$$

Since by hypothesis $u^+$ does not degenerate, there follows that $\alpha>0$.

On the other hand, since $v$ is regular, $v$ admits an asymptotic development 
at the origin in the form $v(x,t)=\beta x_1^+ + o(|x|+|t|^{1/2})$. Clearly, 
$\beta\ge \alpha$.

Now, let $h$ be the caloric function in $\tilde {\cal O}:=
Q_1^-\cap\{ v>0 \}\cap\{-\mu<t<0\}$ for some small $\mu>0$,  with $h=v-u$
on $\partial_p\tilde{\cal O}$. And, let $g$ be the caloric function
in $\tilde{\cal O}$ with $g=v$ on $\partial_p\tilde{\cal O}$. 
Then, $h=g=0$ in $Q_1^-\cap\partial
\{ v>0 \}\cap\{-\mu<t<0\}$ and $h>0$, $g>0$ in $\tilde{\cal O}$.

Therefore, by \cite{ACS}, there exists $\sigma>0$ such that $h\ge\sigma g$
in $Q_{1/2}^-\cap\{ v>0 \}\cap\{-\frac\mu 2<t<0\}$.

Since $u$ is subcaloric in $Q_1^-$ and $u\le v$ in $Q_1^-$ we deduce that
$v-u\ge\sigma u$ in
$Q_{1/2}^-\cap\{ v>0 \}\cap\{-\frac\mu 2<t<0\}$. In particular
$\beta-\alpha\ge\sigma\alpha>0$.


The theorem will be finished if we show that $\alpha=\sqrt{2M(0,0)}$.

\medskip

\underline{Case 1:} $u^{\ve_k}$ verifies (\ref{umenordelta}).
\medskip

As in Theorem \ref{puntual}, we obtain 
\begin{align*}
&\iint_{\cal D} u_t u_{x_1} \psi = \frac12 \iint_{\cal D} |\nabla u|^2 
\psi_{x_1} - \iint_{\cal D} u_{x_1} \nabla u \nabla \psi \\
&\quad+ 
 \iint_{{\cal D}\cap \{u>0\}}M(x,t) \psi_{x_1}+ \iint_{\cald\cap\{u>0\}}
(w_0)_{x_1}\left(\int_{-w_0}^1f(s)ds\right)\psi
\end{align*}
for every test function $\psi$. Then, taking $\psi^{\lambda}(x,t) = 
\lambda \psi(\frac{x}{\lambda},\frac{t}{\lambda^2})$ and changing variables, 
we get 
\begin{align*}
&\iint_{\cal D} (u_\lambda)_t (u_\lambda)_{x_1} \psi = 
\frac12 \iint_{\cal D} |\nabla u_\lambda|^2 \psi_{x_1} - 
\iint_{\cal D} (u_\lambda)_{x_1} \nabla u_\lambda \nabla \psi \\
&\quad+ \iint_{{\cal D}\cap \{u_\lambda>0\}}M(\lambda x,\lambda^2 t) 
\psi_{x_1} +\iint_{\cald\cap\{u>0\}}(w_0)_{x_1}\left(\int_{-w_0}^1f(s)ds\right)
\psi^\lambda
\end{align*}

By Lemma \ref{lim2}, we get (for some sequence $\lambda_k\to0$)
$$0 = -\frac12 \alpha^2 \iint_{{\cal D}\cap\{x_1>0\}} \psi_{x_1} + 
\lim_{k\to\infty}\iint_{{\cal D}\cap\{u_{\lambda_k}>0\}}M(\lambda_k x,
\lambda_k^2 t)\psi_{x_1}.$$

We want to check that $\car_{\{ u_{\lambda_k}>0 \}}\to\car_{\{ x_1>0 \}}$ 
a.e. or, equivalently, 
\begin{enumerate}
\item \label{a}
$\{ x_1>0 \}\subset \cup_{n=1}^{\infty}\cap_{k\ge n}\{ u_{\lambda_k}>0 \} 
= \liminf \{ u_{\lambda_k}>0 \}$ a.e.
\item \label{b}
$\cap_{n=1}^{\infty}\cup_{k\ge n}\{ u_{\lambda_k}>0 \} = 
\limsup \{ u_{\lambda_k}>0 \} \subset \{ x_1>0 \}$ a.e.
\end{enumerate}


Let us see (\ref{a}). If $x_1>0$, we get that $\alpha x_1>0$ and since 
$u_{\lambda_k}(x,t)\to\alpha x_1$ it follows that $u_{\lambda_k}(x,t)>0\quad 
\forall k\ge k_0$.

Let us see (\ref{b}). If exists $k_j\to\infty$ with 
$u_{\lambda_{k_j}}(x,t)>0$ then it must be $x_1\ge0$, because if $x_1<0$, 
we have that $v_{\lambda_{k_j}}(x,t)=0$ for $j\ge j_0$ (because as $v$ is regular, 
$\{ v_{\lambda_k}>0 \}\to \{ x_1>0 \}$). Since $u_{\lambda_{k_j}}\le 
v_{\lambda_{k_j}}$ we get a contradiction.

\medskip

Therefore, 
$$0 = -\frac12\alpha^2\iint_{{\cal D}\cap\{x_1>0\}}\psi_{x_1} + 
M(0,0)\iint_{{\cal D}\cap\{x_1>0\}}\psi_{x_1}.$$
So that, 
$$0=\int_{{\cal D}\cap\{x_1=0\}}\left(\frac12\alpha^2 - M(0,0)\right)\psi dx'dt.$$

Since $\psi$ is arbitrary, $\frac12\alpha^2 = M(0,0)$, so that, 
$$\alpha = \sqrt{2M(0,0)}$$
and the proof is finished

\medskip

\underline{Case 2:} $u^{\ve_k}_t\le 0$
\medskip

We already now that, if we consider $u_{\lambda}(x,t) = \frac{1}{\lambda}
u(\lambda x,\lambda^2 t)$, then it follows that
$$u_{\lambda}(x,t)\to u_0(x,t)\equiv \alpha x_1^+$$
uniformly on compact subsets of $\R^{N+1}$.

As before
\begin{align*}
&\iint_{\cal D} u^{\ve_k}_t u^{\ve_k}_{x_1}\psi = \frac12 
\iint_{\cal D}|\nabla u^{\ve_k}|^2\psi_{x_1} - \iint_{\cal D} u^{\ve_k}_{x_1}
\nabla u^{\ve_k}\nabla\psi
 + \iint_{\cal D}{{\cal B}}_{\ve_k}(u^{\ve_k},x,t)\psi_{x_1}\\
& \quad+ \iint_{\cal D} w^{\ve_k}_{x_1}\left(\int_{-w_0\ve_k}^{u^{\ve_k}} 
f_{\ve_k}(s) ds\right)\psi+\iint_{\cald}(w_0)_{x_1}\big(\frac{\wepk}{\epk}-w_0\big)
f(-w_0)\psi
\end{align*}

Now, as in the previous case, if we consider first $\psi^{\lambda}(x,t) = 
\lambda\psi(\frac{x}{\lambda},\frac{t}{\lambda^2})$ and change variables, 
we obtain
\begin{equation}
\label{velam}
\begin{aligned}
&\iint (u^{\ve_k}_{\lambda})_t (u^{\ve_k}_{\lambda})_{x_1}\psi = 
\frac12 \iint|\nabla u^{\ve_k}_{\lambda}|^2\psi_{x_1} - 
\iint (u^{\ve_k}_{\lambda})_{x_1}\nabla u^{\ve_k}_{\lambda}
\nabla\psi\\
&\quad + \iint{{\cal B}}^\lambda_{\ve_k/\lambda}(u^{\ve_k}_{\lambda}, x,
 t)\psi_{x_1}
 + \iint_{\cald}\frac{ (w^{\ve_k})_{x_1}}{\epk}\left(\int_{-w_0}
^{\frac{u^\ve_k}{\epk}} 
f(s) ds\right)\psi^\lambda\\
&\quad+\iint_{\cald}(w_0)_{x_1}\big(\frac{\wepk}{\epk}-w_0\big)f(-w_0)\psi^\lambda
\end{aligned}
\end{equation}
where $\bep^\lambda(u,x,t)=\int_{-w_0(\lambda x,\lambda^2 t)\ve}^{u}
(s+{w^\ve}(x,t))f_{\ve}(s)ds$.
\medskip
We want to pass to the limit as both $\ve_k$ and $\lambda$ go to zero.

Using Lemma \ref{lim2}, we see that for every sequence
$\lambda_n\to0$ there exists a sequence $k_n\to\infty$ such that $\delta_n:=
\ve_{k_n}/\lambda_n\to 0$ and $u^{\delta_n}:=(u^{\ve_{k_n}})_{\lambda_n}\to
u_0$ uniformly on compact sets of $\R^{N+1}$. By Proposition \ref{lim}  
we see that 
we can pass to the limit in the first three terms of (\ref{velam})
(with $\ve=\ve_{k_n}$ and $\lambda=\lambda_n$).
\medskip

Let us study the limit of
${{\cal B}}^{\lambda_n}_{\delta_n}(u^{\delta_n}(x,t), x,t)$.

It is easy to see that in $\{ x_1>0 \}$, 
${{\cal B}}^{\lambda_n}_{\delta_n}(u^{\delta_n}(x,t),x, t)\to M(0,0)$
uniformly on compact sets. Now, let $K\subset\{ x_1<0\}$ be compact. 
We will show that 
$$\nabla({{\cal B}}^{\lambda_n}_{\delta_n}(u^{\delta_n}(x,t), x, t))\to
0\quad \mbox{in } L^1(K)$$
In fact,
\begin{align*}
&\nabla({\cal B}^{\lambda_n}_{\delta_n}(u^{\delta_n}(x,t), x,t))=
v^{\delta_n}f_{\delta_n}(u^{\delta_n})\nabla u^{\delta_n}\\
&\quad+\lambda_n\nabla w_0(\lambda_nx,\lambda_n^2t) \big(\frac{w^{\delta_n}}
{\delta_n}(x,t)-w_0(\lambda_nx,\lambda_n^2t)\big)f(-w_0(\lambda_n x,\lambda_n^2t))\\
&\quad\qquad+\frac{\nabla w^{\delta_n}}{\delta_n}
\int_{-w_0(\lambda_n x,\lambda_n^2t)}
^{\frac{u^{\delta_n}}{\delta_n}}f(s) ds
\end{align*}
Since $v^{\delta_n}f_{\delta_n}(u^{\delta_n})\to 0$ as measures in $K$ and is
nonnegative, we deduce that the convergence takes place in $L^1(K)$. On
the other hand, $\nabla u^{\delta_n}$ is uniformly bounded. Therefore,
the first term goes to zero in $L^1(K)$.

In order to see that the second and third terms go to zero uniformly
in $K$ we only need to observe that
$$
\frac{u^{\delta_n}}{\delta_n}(x,t)=\frac{\uepkn}{\epkn}(\lambda_nx,\lambda_n^2t)
$$
and a similar formula holds for $\frac{w^{\delta_n}}{\delta_n}$. So that
\begin{align*}
&\big|\frac{w^{\delta_n}}{\delta_n}(x,t)-w_0(\lambda_nx,\lambda_n^2t)|\to0
\quad\mbox{uniformly on compact sets of }\R^{N+1},\\
&\frac{u^{\delta_n}}{\delta_n}\ge-\frac{w^{\delta_n}}{\delta_n}\ge-C,\\
&\frac{|\nabla w^{\delta_n}|}{\delta_n}(x,t)=
\lambda_n\frac{|\nabla\wepkn|}{\epkn}(\lambda_nx,\lambda_n^2t)\to
0\qquad\mbox{uniformly on compact sets of }\R^{N+1}. 
\end{align*}
On the other hand,
$|{\cal B}^{\lambda_n}_{\delta_n}(u^{\delta_n}(x,t),x,t)|\le C_K$, 
so that we have, 
$$
{\cal B}^{\lambda_n}_{\delta_n}(u^{\delta_n}(x,t), x, t)\to \bar M(t) 
\qquad\mbox{weakly in }L^2(K).
$$

Let us now see that, actually,  the convergence takes place in $L^1(K)$.

There holds that 
\begin{align*}
&\frac{\partial}{\partial t}\left({\cal B}^{\lambda_n}_{\delta_n}
(u^{\delta_n}(x,t), x, t)\right) = 
v^{\delta_n}f_{\delta_n}(u^{\delta_n})(u^{\delta_n})_t + 
\frac{\partial}{\partial t}{\cal B}^{\lambda_n}_{\delta_n}
(u^{\delta_n},x, t)\\
&\qquad\quad\le \frac{\partial}{\partial t}{\cal B}^{\lambda_n}_{\delta_n}
(u^{\delta_n}, x, t)
\le C_K\quad\mbox{ in } K.
\end{align*}
On the other hand, for every $(x_0,t_0)\in K,\mbox{ and }\ Q_\tau(x_0,t_0)\subset\{x_1<0\}$
\begin{align*}
\iint_{Q_\tau(x_0,t_0)}&\frac{\partial}{\partial t}\left(
{\cal B}^{\lambda_n}_{\delta_n}(u^{\delta_n}(x,t), x, t)\right) =
\int_{B_\tau(x_0)}{\cal B}^{\lambda_n}_{\delta_n}
(u^{\delta_n}(x,t_0+\tau^2), x,(t_0+\tau^2))\, dx\\
& \qquad- \int_{B_\tau(x_0)}{\cal B}^{\lambda_n}_{\delta_n}
(u^{\delta_n}(x,t_0-\tau^2),x, (t_0-\tau^2))\, dx\\
&\quad\qquad\qquad \ge\ -C_\tau
\end{align*}
since $|{\cal B}^{\lambda_n}_{\delta_n}(u^{\delta_n}(x,t), x, t)|
\le  C_K$ 
for every compact set $K$.

Therefore there exists $C_K>0$ such that $\| {\cal B}^{\lambda_n}_{\delta_n}
(u^{\delta_n}(x,t), x,t)\|_{W^{1,1}(K)}\le C_K$. Hence the convergence takes
place in $L^1(K)$ (for a subsequence).

Now arguing  as in Lemma \ref{Bepsilon}, we get that $\bar M(t)=0$ or 
$\bar M(t)=M(0,0)$.

We can now take the limit in (\ref{velam}) for the sequences $\ve_{k_n}$ and
$\lambda_n$ and  we
obtain
$$0 = -\frac12 \alpha^2\iint_{{\cal D}\cap\{x_1>0\}} \psi_{x_1} + 
M(0,0)\iint_{{\cal D}\cap\{x_1>0\}} \psi_{x_1} + 
\iint_{{\cal D}\cap\{x_1<0\}} \bar M(t) \psi_{x_1}.$$
So that,
$$0 = \int_{{\cal D}\cap\{x_1=0\}}\left(\frac12 \alpha^2 - M(0,0) - \bar M(t)\right)
\psi\, dx' dt.$$

Since $\psi$ is arbitrary we get $\frac12 \alpha^2 = M(0,0) -\bar M(t)$. So that, 
in particular, $\bar M(t)$ is constant and then we have that 
$\bar M(t)\equiv 0$ or $\bar M(t)\equiv M(0,0)$. Since $\alpha>0$ we deduce that $M(t)\equiv 0$ 
and 
$$\alpha = \sqrt{2M(0,0)}.$$
The proof is finished.
\end{pf}





Now we prove a 
proposition that says that, under suitable assumptions, $u^+$ does 
not degenerate on the free boundary. The proof is similar to theorem 6.3 in 
\cite{CLW1}, where the nondegeneracy of $u^+$ was proved in the strictly two 
phase case. Here we assume, instead of $(\ref{umenordelta})$ the somewhat
stronger condition that for every $K\subset\cald$ compact, there exist $0<
\eta<1$
and $\ve_0>0$ such that for every $0<\ve\le\ve_0$
\begin{equation}\label{H2}
\frac{u^\ve}{\ve}\le \eta\qquad\mbox{in }K\cap{\{u\equiv0\}^\circ}.
\end{equation}


\begin{prop} 
\label{degeneracy}
Let $u=\lim u^{\ve_k}$, where $(u^{\ve_k},v^{\ve_k})$ are uniformly bounded
solutions to $(1.1)$ satisfying $(\ref{ochica})$ with $w_0>-1$, such that
$v^{\ve_k}\ge0$
and the functions $u^{\ve_k}$ satisfy $(\ref{H2})$.
Let $(x_0,t_0)\in\partial\{ u>0 \}$. 
 
Let us assume that there exists $\nu\in \R^N$, with $|\nu|=1$ such that
$$\liminf_{r\to 0^+}\frac{|\{u>0\}\cap\{\langle x-x_0,\nu\rangle>0\}\cap
Q^-_r(x_0,t_0)|}{|Q^-_r(x_0,t_0)|}>\alpha_1$$
and
$$\liminf_{r\to 0^+}\frac{|\{ u=0 \}^\circ\cap\{\langle x-x_0,\nu\rangle<0\}\cap
Q^-_r(x_0,t_0)|}{|Q^-_r(x_0,t_0)|}>\alpha_2$$
with $\alpha_1+\alpha_2>\frac12$, then there exists a constant $C>0$ and 
$r_0>0$ which depends on $N$ and $f$ such that, if $0<r\le r_0$, 
$$\sup_{\partial_p Q_r^-(x_0,t_0)} u\ge C r.$$
\end{prop}

\begin{pf}
Without loss of generality, we may assume that $(x_0,t_0)=(0,0)$ and that  
$\nu=e_1=(1,0,...,0)$. 

We will note $Q_r^-=Q_r^-(0,0)$ and
$$
(u^\ve)_r(x,t)=\frac{1}{r}u^\ve(rx,r^2t),\quad (v^\ve)_r(x,t)=\frac{1}{r}
v^\ve(rx,r^2t),\quad
u_r(x,t)=\frac{1}{r}u(rx,r^2t).$$

Let us see that there exists $r_0>0$ and a constant $c$ such that if
$r<r_0$ and $\ve<\ve_0=\ve_0(r)$, then
$$\iint_{Q_1^-}(v^\ve)_r f_{\ve/r}((u^\ve)_r) dx\ge c.$$


Following \cite{CLW1}, there exists $\gamma>0$ small such that, for some 
$\lambda>0$, 
$$\frac{|\{u_r>\gamma\}\cap\{ x_1>0 \}\cap Q^-_1|}{|Q^-_1|} + 
\frac{|\{ u_r=0 \}^\circ\cap\{ x_1<0 \}\cap Q^-_1|}{|Q^-_1|}\ge \frac12 + \lambda$$

Let us now define
$$A_r=\{ u_r>\gamma \}\cap\{ x_1>0 \}\cap Q_1^-,\quad
B_r=\{ u_r=0 \}^\circ\cap\{ x_1<0 \}\cap Q_1^-$$
$$\mbox{and } -B_r=\{(x_1,x',t)/\ (-x_1,x',t)\in B_r\}.$$
Then, we have
$$|A_r\cap (-B_r)|\ge \lambda |Q_1^-| = \tilde{\lambda}.$$
Once again, following \cite{CLW1} we have for $0<\rho<1$ fixed, that there  
exists $0<x_1^r<1$ such that
$$|\Lambda_r|=|\{ (x',t)/\ (x_1^r,x',t)\in A_r\cap (-B_r)\}|>
\rho\tilde{\lambda}$$

Let $\eta>0$ be the constant in $(\ref{H2})$ in $Q_1(0,0)$, let
$0<\delta'<\delta$, $0<b<b'<1$ be such that
$$
\eta< -w_0(0,0)+\delta<b.
$$
Let $\kappa>0$ be such that
$$f(s)>\kappa>0\quad \mbox{for } s\le b'.$$
Then, for $(x',t)\in\Lambda_r$ we have
$$\frac{1}{\ve/r}(u^\ve)_r(x_1^r,x',t)>\frac{\gamma}{2(\ve/r)}>b,\quad
\frac{1}{\ve/r}(u^\ve)_r(-x_1^r,x',t)< -w_0(0,0)+\delta$$
if $\ve_k<\ve_1=\ve_1(r)$ is small.
 So that, for
every $(x',t)\in \Lambda_r$ there exists $\tilde x_1^r\in(-1,1)$ such that
$-w_0(0,0)
+\delta\le\frac1{\ve/r}(u^\ve)_r(\tilde x_1^r,x',t)\le b$.

Now, by the uniform Lipschitz regularity of $(u^\ve)_r$ and $(v^\ve)_r$,  and
$(\ref{ochica})$,
we have that for $\ve\le\ve_0(\le\ve_1)$ and $r\le r_0$,
$$
\frac{(u^\ve)_r}{\ve/r}(x_1,x',t)\le b'\quad
\mbox{and}\quad\frac{(v^\ve)_r}{\ve/r}
(x_1,x',t)\ge \delta'
\quad \mbox{if }
|x_1-\tilde x_1^r|<C\frac{\ve}{r}
$$
where $C$ depends on $\delta,\delta',b,b'$, on the Lipschitz constant of 
$u^\ve$ and $v^\ve$ in $Q_1^-$  and $r_0$ depends only $w_0$.

Finally we have
\begin{align*}
\iint_{Q_1^-}(v^\ve)_r &f_{\ve/r}((u^\ve)_r)=
\iint_{Q_1^-}\frac{(v^\ve)_r}{\ve/r}\frac1{\ve/r}f\big(
\frac{(u^\ve)_r}{\ve/r}\big)\\
&\ge \delta'\frac{\kappa}{\ve/r}
\left|\left\{(x,t)\in Q_1^-/\ \frac{(v^\ve)_r}{\ve/r}\ge\delta'\mbox{ and }
f\big(\frac{(u^\ve)_r}{\ve/r}\big)\ge \kappa\right\}\right|\\
&\qquad\qquad\qquad\ge\delta'\frac{\kappa}{\ve/r}|\Lambda_r|2C\frac{\ve}{r}\ge
2C\delta'\kappa\rho\tilde{\delta}\equiv c.
\end{align*}

The rest of the proof follows as in \cite{CLW1}
\end{pf}

\begin{rem}
\label{otrodegen}
 Proposition \ref{degeneracy} remains true if we change the 
hypothesis that $\uepk$ satisfies $(\ref{H2})$ by 
\begin{equation}
\label{H3}
\frac{u^{\ve_k}}{\ve_k}\to -w_0\qquad a.e. \ \ \{u\equiv0\}^\circ.
\end{equation}
In fact, as in the proof of Prop. \ref{degeneracy} we consider for each
$0<r<1$ the sets $A_r$ and $B_r$. So that, for some $0<\lambda<1$
$$
|A_r\cap(-B_r)|\ge\lambda |Q_1^-|.
$$
Since $B_r\subset\{u_r\equiv0\}^\circ$, there holds that
$$
\frac{(u^\ve)_r}{\ve/r}(-x_1,x',t)\to -w_0(-r{x_1},rx',r^2t)\qquad a.e. \ \ A_r\cap(-B_r).
$$
Let $0<\mu<1$. There exists $C_r\subset \big(A_r\cap(-B_r)\big)$ such that
$|C_r|=\mu|A_r\cap(-B_r)|$ and 
$$
\frac{(u^\ve)_r}{\ve/r}(-x_1,x',t)\to -w_0(-rx_1,rx',r^2t)\quad\mbox{uniformly in }C_r.
$$
Let $\delta>0$. There exists $\ve_1=\ve_1(r)$ such that
$$
\frac{(u^\ve)_r}{\ve/r}(-x_1,x',t)\le -w_0(-rx_1,rx',r^2t)+\frac\delta2\le
-w_0(0,0)+\delta\quad\mbox{in }C_r
$$
if $\ve<\ve_1$ and $r<r_0= r_0(\delta)$.
Now, the proof follows as in Prop. \ref{degeneracy} by taking $\tilde\lambda=
\mu\lambda|Q_1^-|$ and 
$$
\Lambda_r:=\{(x',t)\,/\,(x_1^r,x',t)\in C_r\}.
$$
\end{rem}


\begin{rem}
\label{otrodegen2}
Proposition \ref{degeneracy} remains true if we change
condition $(\ref{H2})$ by condition $(\ref{umenordelta})$. In fact,
as in the proof of Theorem \ref{puntual} we see that condition 
$(\ref{umenordelta})$ implies that 
$$
\bepk(\uepk, x,t)\to0\qquad L^1_{loc}(\{u\equiv0\}^\circ).
$$
As in Theorem \ref{puntual} we deduce that $\uepk$ satisfies $(\ref{H3})$.

\end{rem}

\bigskip


Using  Remark \ref{otrodegen}, Remark \ref{otrodegen2} and Theorem \ref{viscosa}
 we get the following
Corollaries.




\begin{cor}
\label{cor1}
Let $u=\lim u^{\ve_k}$ where $(u^{\ve_k},v^{\ve_k})$ are unifomly bounded
solutions to $(1.1)$ in a domain $\cald\subset\R^{N+1}$ with $v^{\ve_k}\ge0$,
which verify $(\ref{ochica})$ with $w_0>-1$ and such that 
$u^{\ve_k}$ satisfies $(\ref{umenordelta})$.
If the free
boundary $\cald\cap\partial\{u>0\}$ is given by $x_1 = g(x',t)$ with
$g\in Lip(1,1/2)$, then, $u$ is a 
viscosity solution of the free boundary problem $(\ref{FBP})$.
\end{cor}




\begin{cor}
\label{cor0}
Let $u=\lim u^{\ve_k}$ where $(u^{\ve_k},v^{\ve_k})$ are unifomly bounded
solutions to $(1.1)$ in a domain $\cald\subset\R^{N+1}$ with $v^{\ve_k}\ge0$,
which verify $(\ref{ochica})$ with $w_0>-1$ and such that $\uepk$ 
satisfies $(\ref{H3})$ and $u^{\ve_k}_t\le 0$.
If, for every $(x_0,t_0)\in \cald\cap\partial\{u>0\}$,  
$\{x\in \R^N\,/\, (x,t_0)\in \cald\cap\{u>0\}\}$ is given by 
$x_1>\Phi(x')$ with
$\Phi$, Lipschitz continuous  then, $u$ is a 
viscosity solution of the free boundary problem $(\ref{FBP})$.
\end{cor}
\begin{pf} We only need to see that $u$ does not degenerate at points
of the free boundary which are regular from the zero side. Let $(x_0,t_0)$ be
any such point. We see that
we can apply Remark \ref{otrodegen} at that point. In fact, since $\uepk_t
\le0$, $u$ is decreasing in time. Therefore,
$$
\{(x,t)\,/\,x_1>\Phi(x)\,,\,t\le t_0\}\subset\{u>0\}
$$  
and the parabolic density of this set is positive.
\end{pf}



In particular, Corollary \ref{cor0} can be applied to solutions of (1.1) with
$u_0^\ve$
constructed as in \cite{CV} and $v_0^\ve$ a small perturbation of $u_0^\ve$.


\begin{cor}
\label{cor4}
Let $u_0\in C(\R^N)\cap C^2(\overline{\{u_0>0\}})$ be such that
 $||u_0||_{C^2(\overline{\{u_0>0\}})}<\infty$, $ \Delta u_0\le0$ and
$(u_0)_{x_1}-\lambda|\nabla u_0|\ge 0$ in $\{u_0>0\}$ with $\lambda>0$.
Assume, moreover that   $0<a_2\le |\nabla u_0|\le a_1<\sqrt{2M_0}$ in
a neighborhood of the free boundary: $\{x\in \{u_0>0\}\,/\, dist(x,
\{u_0=0\})\le \gamma\}$, and $M_0=\int_0^1sf(s)$. Then, there exists a 
sequence $(u_0^\ve,v_0^\ve)\in \big(C^1(\R^N)\big)^2$
with $u_0^\ve\to u_0$ uniformly in $\R^N$ (so that $u_0^\ve$ are uniformly
bounded) and, moreover, they satisfy
\begin{equation}
\label{u0ep}
\begin{aligned}
&1)\ \Delta u_0^\ve-v_0^\ve f_\ve(u_0^\ve)\le0\\
&2)\ (u_0^\ve)_{x_1}-\lambda|\nabla u_0^\ve|\ge 0\\
&3)\ \frac{v_0^\ve-u_0^\ve}{\ve}\to w_0\ \ \mbox{uniformly on compact sets, with
}\ \ w_0>-1.
\end{aligned}
\end{equation} 
 $w_0\in \R$ is any constant  such that $w_0\ge-\eta$ with $\eta>0$ small
enough.


Let $(u^\ve,v^\ve)$ be the solution to $(\ref{1.1})$ with initial datum 
$(u_0^\ve,v_0^\ve)$ (so that, in particular, $u^\ve$ and $v^\ve$ are
uniformly bounded).  For every sequence $\epj\to 0$ there exists
a subsequence $\epjk$ such that there exists
$$
u=\lim_{k\to\infty}\uepjk
$$
and $u$ is a viscosity solution to the free boundary problem $(\ref{FBP})$.

\end{cor}

\begin{pf}  Let $u_0^\ve$ be the approximations constructed in \cite{CV}.
The approximations are constructed in the following way. First we extend
$u_0$ to a neighborhood of $\{u_0>0\}$: $S:=\{x\in \R^N\,/\,\mbox{dist}\,(x,
\{u_0>0\})\le\gamma\}$ in such a way that $||u_0||_{C^2(S)}<\infty$. 
For $\ve$ small enough we define
$$
u_0^\ve(x)=\ve F\left(\frac1{\sqrt{2M_0}}\big(1-\frac{u_0(x)}{\ve}\big)\right)
\qquad\mbox{in }\ \{-C\ve\le u_0\le \ve\}.
$$
where $F\in C^2(\R)$ is such that
$$
F''\le (1+\delta)F f(F)+\alpha F',\ \ F(0)=1,\ \ F'(0)=-\sqrt{2M_0}.
$$
Here $\delta>0,\ \alpha>0$ are such that $F$ has a strict minimum at
a finite point $\bar s$ such that $\bar s\sqrt{2M_0}>1$. ($\bar s\to+
\infty$ as $\delta\to 0$), and $F$ is decreasing for $s<\bar s$. 

The constant $C$ is taken as $C=\bar s\sqrt{2M_0}-1$.

We define 
\begin{align*}
u^\ve_0&=u_0\ \,\qquad\mbox{in }\ \{u_0>\ve\}\\
u_0^\ve&=\ve F(\bar s)\quad\mbox{in }\ \R^N\setminus\{u_0>-C\ve\}.
\end{align*}


As in \cite{CV}, we see that  $u_0^\ve\in C^1(\R^N)$.

Let $w_0\in\R$ be such that $w_0\ge-\eta>-F(\bar s) $  with
$\eta>0$ to
be fixed  later and let
$$
v_0^\ve=u_0^\ve+\ve w_0.
$$
Then, $v_0^\ve\ge0$. 
It is immediate to verify that (\ref{u0ep}) 1) is satisfied
in $\{u_0>\ve\}$ and $\R^N\setminus\{u_0>-C\ve\}$. Let us see that it is
satisfied in $\{-C\ve\le u_0\le \ve\}$. In fact,
\begin{align*}
\Delta u_0^\ve&-v_0^\ve f_\ve(u_0^\ve)=\frac1{2M_0\ve}F''|\nabla u_0|^2-
\frac1{\sqrt{2M_0}}F'\Delta u_0
-\frac1\ve F f(F)-\frac{w_0}\ve f(F)\\
&\le \frac{1+\delta}{2M_0\ve}
F f(F)|\nabla u_0|^2
+\frac\alpha{2M_0\ve}F'|\nabla u_0|^2-\frac a{\sqrt{2M_0}} F'-
\frac1\ve F f(F) -\frac{w_0}\ve f(F)
\end{align*}
where $a>0$ is such that $|\Delta u_0|\le a$.


Let $0<\mu<1$  be such that $a_1\le (1-\mu)^{1/2}A\sqrt{2M_0}$ with $0<A<1$,
and let $\delta$ in the definition of $F$ be such that $(1+\delta)A^2\le 1$.
Then, if $\ve$ is small enough so that ${\alpha a^2_2}/\sqrt{2M_0}>{a\ve}$
there holds that
\begin{align*}
\Delta u_0^\ve&-v_0^\ve f_\ve(u_0^\ve)\le \frac1\ve\left[[(1+\delta)(1-\mu)A^2
-1]F f(F)+\big(\frac{\alpha a^2_2}{2M_0}-\frac{a\ve}{\sqrt{2M_0}}\big) F'-w_0 f(F)
\right]\\
&\le \frac1\ve[-\mu F-w_0]f(F)\le \frac1\ve[-\mu F(\bar s)-w_0]f(F)\le0
\end{align*}
 if $\eta=\mu F(\bar s)$.

Clearly, (\ref{u0ep}) 3)  holds. Let us see that (\ref{u0ep}) 2) also
holds. We only need to verify this property in the set $\{-C\ve< u_0<\ve\}$
and this is clear from the fact that
$$
\nabla u_0^\ve=-\frac1{\sqrt{2M_0}}F'\left(\frac1{\sqrt{2M_0}}\big({1-
\frac{u_0}\ve}\big)\right)\nabla u_0.
$$


Now, by the results of Section 2, for every sequence
$\epj\to0$ there exists a subsequence and a continuous function $u$ such
that $\uepjk\to u$ uniformly on compact subsets of $\R^N\times(0,\infty)$.


On the other hand,
 $u^\ve_t$ is a solution to the following equation
$$
\Delta U-U_t=\beta_\ve'(u)U.
$$
Here $\beta_\ve  (s)=sf_\ve(s)$. Since, for $\ve$ small enough $u^\ve_t(x,0)\le0$
we conclude that
\begin{equation}
\label{ut}
u^\ve_t\le0\qquad\mbox{in }\R^N\times(0,\infty).
\end{equation}

In a similar way we see that $u_{x_1}-\lambda u_{x_i}\ge0$ for every $i$. So
that
\begin{equation}
\label{ux1}
u^\ve_{x_1}-\frac\lambda N|\nabla u^\ve|\ge0\qquad\mbox{in }\R^N\times(0,\infty).
\end{equation}

Clearly (\ref{ut}) and (\ref{ux1}) imply that
$$
u_t\le0\qquad\mbox{and }\quad u_{x_1}-\frac\lambda N|\nabla u|\ge0\qquad
\mbox{in }\{u>0\}.
$$

In particular, the free boundary is Lipschitz in space.

So that, in order to apply Corollary \ref{cor0} we only need to verify that
$\uepk$ satisfies (\ref{H3}).
On one hand, given $K\subset\{u_0\equiv0\}^\circ$ compact, there exists
$\ve_0$ such that for $\ve<\ve_0$
$$
{\cal B}_\ve(u^\ve_0,x,0)=\int_{-w_0}^{\frac{u_0^\ve}\ve(x)}(s+w_0)
f(s)=\int_{-w_0}^{F(\bar s)}(s+w_0)f(s).
$$
On the other hand,
$$
\frac\partial{\partial t}\big({\cal B}_\ve(u^\ve,x,t)\big)=
v^\ve f_\ve(u^\ve)u^\ve_t\le0.
$$
Therefore,
$$
{\cal B}_\ve(u^\ve,x,t)\le \int_{-w_0}^{F(\bar s)}(s+w_0)f(s)\qquad\mbox{for }
x \mbox{ in }K\,,\, t>0.
$$

As in the proof of Theorem \ref{viscosa} we see that, since $u^\ve_t\le0$,
there holds that ${\cal B}_\ve(u^\ve,x,t)\to \bar M(x,t)$ in $L^1_{loc}(\{u\equiv
0\}^\circ)$ and, for almost every $(x,t)$ we either have $\bar M(x,t)=0$ or $\bar M(x,t)=M
=\int_{-w_0}^1(s+w_0)f(s)$.
Since
$$
\nabla \big(B_\ve(u^\ve,x,t)\big)=v^\ve f_\ve(u^\ve)\nabla u^\ve\to 0\quad\mbox
{in }L^1_{loc}(\{u\equiv0\}^\circ)
$$
there holds that $\bar M(x,t)=\bar M(t)$ in $\{u\equiv0\}^\circ$. Therefore,
$$
\bar M(t)\le \int_{-w_0}^{F(\bar s)}(s+w_0)f(s)\quad a.e.\ \ \{u\equiv0\}^\circ.
$$

Since $F(\bar s)<1$, there holds that $\bar M(t)\equiv0$. 

Thus, for every sequence $\epk\to 0$
$$
\int_{-w_0}^{\frac{\uepk}{\epk}}(s+w_0)f(s)\to 0\qquad a.e.\ \  \{u\equiv0\}^\circ
$$
and we deduce that $\uepk$ satisfies (\ref{H3}).
\end{pf}






Combining the regularity results for viscosity solutions of \cite{F},
Corollary \ref{cor1} and Corollary \ref{cor0} we have
the following regularity result for limit functions.

\begin{cor}
\label{cor3}
Let $u$ as in Corollary \ref{cor1} or Corollary \ref{cor0}.
If, moreover,  the free
boundary $\cald\cap\partial\{u>0\}$ is given by $x_1 = g(x',t)$ with $g$
 Lipschitz continuous,
then, $u$ is a classical solution of the free boundary problem $(\ref{FBP})$.
\end{cor}



\begin{ack}
The first author wants to thank J. D. Rossi for his constant support and 
several interesting discussions.
\end{ack}

\end{section}








\begin{thebibliography}{[xxx]}

\bibitem{ACS} I. Athanasopoulos, L.A. Caffarelli, S. Salsa, {\it Caloric 
functions in Lipschitz domains and the regularity of solutions to phase 
transitions problems}, Ann. Math. 143(3) (1996), 413--434.

\bibitem{BCN} H. Berestycki, L.A. Caffarelli, L. Nirenberg, 
{\it Uniform estimates for regularization of free boundary problems}, Annal. 
and Partial Diff. Eq., Cora Sadosky eds., Lecture Notes in Pure and Applied 
Math., V. 122, Marcel Dekker, 1988.

\bibitem{BL} J.D. Buckmaster, G.S.S. Ludford, Theory of laminar flames, 
Cambridge Univ. Press, Cambridge, 1982.

\bibitem{C} L.A. Caffarelli, {\it Uniform Lipschitz regularity of a 
singular perturbation problem}, Diff. and Int. Eq., V. 8, N. 7, (1995), 
1585--1590.

\bibitem{CLW1} L.A. Caffarelli, C. Lederman, N. Wolanski, {\it Uniform 
estimates and limits for a two phase parabolic singular perturbation
problem}, Indiana Univ. Math. J. 46 (2) (1997), 453--490.

\bibitem{CLW2} L.A. Caffarelli, C. Lederman, N. Wolanski, {\it Pointwise 
and viscosity solutions for the limit of a two phase parabolic singular 
perturbation problem}, Indiana Univ. Math. J. 46 (3) (1997), 719--740.

\bibitem{CV} L.A. Caffarelli, J.L. V\'azquez, {\it A free boundary 
problem for the heat equation arising in flame propagation}, Trans. Amer. 
Math. Soc. {\bf 347} (1995), 411--441.

\bibitem{F} L. Fornari, {\it Regularity of the solution and of the free 
boundary for free boundary problems arising in combustion theory}, preprint.


\bibitem{LSU} O.A. Ladyzhenskaya, V.A. Solonnikov, N.N. Ural'tseva, 
Linear and Quasilinear Equations of Parabolic Type, Transl. Math. 
Monographs, Vol. 23, Amer. Math. Soc., Providence, R.I., 1968.

\bibitem{LVW} C. Lederman, J. L. Vazquez, N. Wolanski, {\it Uniqueness
of solution in a free boundary problem from combustion}, to appear in Trans. AMS.

\bibitem{LVW2} C. Lederman, J. L. Vazquez, N. Wolanski, {\it Uniqueness
of solution to a two phase free boundary problem} preprint.

\bibitem{LW} C. Lederman, N. Wolanski, {\it Viscosity solutions and 
regularity of the free boundary for the limit of an elliptic two phase 
singular perturbation problem},  Analli della Scuola Normale
Sup. Pisa. Serie IV, Vol. XXVII. Fasc. 2, (1998), 253--288.

\bibitem{ZF} Ya.B. Zeldovich, D.A. Frank-Kamenetski,
{\it The theory of thermal propagation of flames},
 Zh. Fiz. Khim. 12,  (1938), 100--105  (in Russian);
English translation in ``Collected Works of  Ya.B.   Zeldovich",
vol. 1, Princeton Univ. Press,  1992.


\end{thebibliography}


\end{document}
